% Standalone source: compile with pdflatex until references stabilize.
% All sections and the bibliography are embedded.
\documentclass[letterpaper,11pt]{article}
\usepackage[utf8]{inputenc}
\usepackage[english]{babel}
\usepackage[margin=1in]{geometry}
\usepackage{amsmath,amsfonts,amssymb,amsthm,mathtools}
\usepackage{microtype}
\usepackage{needspace}
\usepackage{booktabs}
\usepackage[dvipsnames]{xcolor}
\definecolor{linkblue}{HTML}{0072BC}
\definecolor{citegreen}{HTML}{3C7F32}
\usepackage[colorlinks=true,pdfpagemode=UseNone,urlcolor=linkblue,
  linkcolor=linkblue,citecolor=citegreen,pdfstartview=FitH,
  pdftitle={Improved Range Avoidance for Local Circuits},
  pdfauthor={Yan Zhong}]{hyperref}
\usepackage[nameinlink,noabbrev]{cleveref}
\setlength{\emergencystretch}{1.5em}
\setcounter{tocdepth}{2}
\newtheorem{theorem}{Theorem}
\newtheorem{lemma}[theorem]{Lemma}
\newtheorem{corollary}[theorem]{Corollary}
\newtheorem{proposition}[theorem]{Proposition}
\theoremstyle{definition}
\newtheorem{definition}[theorem]{Definition}
\newtheorem{remark}{Remark}
\newcommand{\F}{\mathbb F}
\newcommand{\R}{\mathbb R}
\newcommand{\Q}{\mathbb Q}
\newcommand{\E}{\mathbb E}
\newcommand{\U}{\mathbf U}
\newcommand{\NC}{\mathsf{NC}}
\newcommand{\FP}{\mathsf{FP}}
\newcommand{\NP}{\mathsf{NP}}
\newcommand{\Avoid}{\textnormal{\textsc{Avoid}}}
\newcommand{\norm}[1]{\left\lVert#1\right\rVert}
\newcommand{\abs}[1]{\left|#1\right|}
\DeclareMathOperator{\tr}{tr}
\DeclareMathOperator{\diag}{diag}
\DeclareMathOperator{\supp}{supp}
\DeclareMathOperator{\rank}{rank}
\DeclareMathOperator{\poly}{poly}
\DeclareMathOperator{\sgn}{sgn}
\DeclareMathOperator{\Range}{Range}

\title{Logarithm-Free Range Avoidance for Local Circuits}
\author{Yan Zhong\thanks{Johns Hopkins University,
  \texttt{\href{mailto:yanzhong.cs@gmail.com}{yanzhong.cs@gmail.com}}.}}
\date{September 23, 2026}

\begin{document}
\hypersetup{pageanchor=false}
\begin{titlepage}
\maketitle
\thispagestyle{empty}
\begin{abstract}
The range avoidance problem asks, given a Boolean circuit with more
outputs than inputs, to find a string outside its range. Guruswami,
Lyu, and Yuan gave a deterministic polynomial-time algorithm for
circuits of locality $k\ge3$ when the output length is at least
$c^k n^{(k-1)/2}\log n$. We remove the logarithmic factor for every
fixed locality. For each fixed $k\ge3$, there is a constant $C_k$ such
that range avoidance is solvable in deterministic polynomial time
whenever $m\ge C_k n^{(k-1)/2}$. The result allows arbitrary local
predicates, arbitrary input neighborhoods, and repeated neighborhoods.
In particular, locality three admits a sufficiently large constant
multiple of linear stretch, and locality four admits stretch
$C_4n^{3/2}$.

Our main technical result constructs efficiently evaluable certificates
for cancellation in signed XOR systems with arbitrary supports and
bounded rational coefficients. At density $M\ge C_{d,\epsilon}n^{d/2}$,
the certificate dominates the maximum absolute value of the signed
polynomial for every signing, while its expectation is at most
$\epsilon M/2$ under every sufficiently small-biased distribution on
the original signs. This allows all Fourier layers of a local circuit
to use one common signing.

The proof develops a constructive use of the spectral hypergraph Moore
argument of Schmidhuber and Hastings. Large trace moments expose a small
window through deterministic recovery of heavy Fourier coefficients
of walk matrices. For odd XOR degree, equalizing the total weight of
each original pair preserves the polynomial-to-matrix identity under
pair deletion. A failed local capacity condition yields a dense forest
of pairs, whose edge signs are independent under uniform original
signs. Removing these forests leaves a remainder with a small trace
certificate. We also obtain the tradeoff
$m\ge C_k n(n/\ell)^{(k-3)/2}$ in time
$\poly(\text{input})(\mathrm e n/\ell)^{O_k(\ell)}$ for
$1\le\ell\le n$.

The certificates give polynomial-size one-sided distinguishers that
reject the entire image and accept at least half of uniform strings,
excluding local demi-bits generators at the same stretch. They also
give explicit query problems requiring $\Omega_t(m^{2/(t-1)})$ bits
of storage with $t$ nonadaptive bit probes; three probes require
linear storage. A separate full-degree argument gives a
logarithm-free $\Omega_{t,w}(m^{2/t})$ bound for adaptive probes
to fixed-size words.
\end{abstract}
\vfill
\begin{center}\small Research draft\end{center}
\end{titlepage}
\pdfbookmark[1]{Contents}{contents}
{\small\tableofcontents}
\thispagestyle{empty}
\clearpage
\setcounter{page}{1}
\hypersetup{pageanchor=true}

\section{Introduction}
\label{sec:introduction}

Let $f:\{0,1\}^n\to\{0,1\}^m$ be a Boolean circuit with $m>n$.
The range avoidance problem, denoted by $\Avoid$, asks for a string
$y\notin\Range(f)$. Such a string exists by counting, but the counting
argument does not supply a deterministic efficient construction.
For a fixed integer $k$, the problem $\NC^0_k$-$\Avoid$ restricts
each output coordinate to depend on at most $k$ input variables.
The local functions and their input neighborhoods may differ between
coordinates. We use the explicit local representation: each coordinate
is specified by its neighborhood and its truth table.

Range avoidance connects explicit constructions, pseudorandomness,
and circuit lower bounds. Guruswami, Lyu, and Wang~\cite{GLW22}
developed these connections for low-depth circuits and gave an
efficient algorithm at locality two. Gajulapalli, Golovnev, Nagargoje,
and Saraogi~\cite{GGNS23} established both hardness consequences and
deterministic algorithms for constant locality. Their algorithm applies
when $m\ge n^{k-1}/\log n$. More recently, Guruswami, Lyu, and
Yuan~\cite[Theorem~5]{GLY26} obtained a deterministic polynomial-time
algorithm when
\begin{equation}
 m\ge c^k n^{(k-1)/2}\log n.
 \label{eq:previous-threshold}
\end{equation}
Their proof reduces range avoidance to cancellation in signed XOR
systems and derandomizes the analysis of semirandom XOR
refutation~\cite{GKM22,HKM23}.

The logarithmic factor in Eqn.~\eqref{eq:previous-threshold} is
particularly visible at small locality. At locality three it separates
linear stretch from $n\log n$; at locality four it separates
$n^{3/2}$ from $n^{3/2}\log n$. We show that this factor can be
removed for every fixed locality.

\subsection{Our Results}

\begin{theorem}[Range avoidance]
\label{thm:avoid}
For every fixed integer $k\ge3$, there is a constant $C_k>0$ and a
deterministic algorithm with the following guarantee. Given an explicit
$k$-local map $f:\{0,1\}^n\to\{0,1\}^m$, $n\ge1$, with
\begin{equation}
 m\ge C_k n^{(k-1)/2},
 \label{eq:main-stretch}
\end{equation}
the algorithm finds a string outside $\Range(f)$ in time polynomial in
the input length. The local functions and neighborhoods are arbitrary,
and repeated neighborhoods are allowed.
\end{theorem}

The constants and the polynomial running-time exponent may depend on
$k$. Thus \Cref{thm:avoid} is a fixed-locality result. It gives
$m\ge C_3n$ at locality three and $m\ge C_4n^{3/2}$ at locality four.
It does not give polynomial time for every $m>n$. In particular,
the sufficiently large constant $C_3$ is essential to the statement
proved here.

The technical result is a signing theorem with an expectation guarantee
that remains valid under a common pseudorandom distribution. For
$F\subseteq[n]$ and $x\in\{\pm1\}^n$, write
$\chi_F(x)=\prod_{v\in F}x_v$.

\begin{theorem}[Certificates for signed XOR systems]
\label{thm:signing}
Fix an integer $d\ge2$ and a rational $\epsilon\in(0,1)$.
There is a constant $C_{d,\epsilon}>0$ with the following property.
Given $n\ge d$, labeled sets $F_1,\ldots,F_M\in\binom{[n]}d$, and rational
coefficients $a_i\in[-1,1]$, with
\begin{equation}
 M\ge C_{d,\epsilon}n^{d/2},
 \label{eq:signing-density}
\end{equation}
a deterministic algorithm constructs a nonnegative rational-valued
certificate $\mathcal C:\{\pm1\}^M\to\Q_{\ge0}$ and a rational
$\eta>0$ such that
\begin{align}
 \sup_{x\in\{\pm1\}^n}
 \abs{\sum_{i=1}^M a_i b_i\chi_{F_i}(x)}
 &\le\mathcal C(b) &&\text{for every }b,\label{eq:signing-soundness}\\
 \E_{b\sim\mathcal D}\mathcal C(b)
 &\le\epsilon M/2
 &&\text{for every fully $\eta$-biased distribution }\mathcal D.
 \label{eq:signing-expectation}
\end{align}
The construction, certificate evaluation, and $\eta^{-1}$ are polynomial
in the explicit input size, with constants depending on $d,\epsilon$.
Labels remain distinct when their supports coincide.
\end{theorem}

A fully $\eta$-biased distribution is one under which every nonempty
parity of the sign coordinates has expectation of magnitude at most
$\eta$. The certificate is constructed before selecting the final
signing. Its decomposition may use a fixed explicit biased space, but
Eqn.~\eqref{eq:signing-expectation} holds for every distribution with
the stated bias. This quantifier permits certificates from several
systems to be summed before one common signing is selected. The
statement includes zero coefficients, negative coefficients, and
arbitrary repetitions.

The level of the matrix construction gives a time--stretch tradeoff.
Throughout the paper, $\mathrm e$ denotes the base of the natural
logarithm.

\begin{theorem}[Hierarchy tradeoff]
\label{thm:hierarchy-avoid}
For each fixed $k\ge3$, there are constants $C_k,c_k>0$ such that,
for every integer $1\le\ell\le n$, range avoidance for explicit $k$-local
maps is solvable deterministically when
\begin{equation}
 m\ge C_k n(n/\ell)^{(k-3)/2},
 \label{eq:hierarchy-stretch}
\end{equation}
in time
\begin{equation}
 \poly(\text{input})\,(\mathrm e n/\ell)^{c_k\ell}.
 \label{eq:hierarchy-time}
\end{equation}
\end{theorem}

The polynomial in Eqn.~\eqref{eq:hierarchy-time} can be chosen with
degree depending only on $k$. Taking $\ell$ constant recovers
\Cref{thm:avoid}; increasing $\ell$ reduces the required output length
while increasing the running time. At $\ell$ comparable to $n$,
the algorithm may enumerate all input assignments. No polynomial-time
claim at nearly minimal stretch follows from this endpoint.

For $k\ge4$, taking a sufficiently large but constant level yields
any prescribed positive coefficient $\rho n^{(k-1)/2}$ in polynomial
time, with exponent depending on $\rho$; see \Cref{cor:prefactor}.
This is a constant-factor consequence, not a reduction of the power
of $n$. At $k=3$, increasing the level does not change the density.

\paragraph{Cryptographic and data structure consequences.}
Our proof certifies a large set of nonoutputs, a stronger guarantee
than producing one missing string. In \Cref{thm:onesided} we construct
a polynomial-size deterministic test $D_f$ that rejects every image
point and accepts at least half of uniformly random strings.
It also accepts at least one third of any distribution with a
sufficiently small polynomial full bias. The construction is uniform
in the local description of $f$.

This immediately breaks demi-bits security against $\NP/\poly$
in the stretch regime of the main theorem. For example, every
five-local Goldreich generator at $m=\rho n^2$, with fixed $\rho>0$,
has such a one-sided attack. This does not exclude its smaller-stretch
instantiations or dense constant-degree polynomial generators.
We use the demi-bits definition and candidate discussion
of~\cite{RWZ26}. Li, Ren, and Zhong~\cite{LRZ26} give a
locality-preserving projection reduction from demi-bits to hard
range-avoidance instances. Applying their projection mechanism to
our algorithm yields polynomial-size nondeterministic certificates
for a $1-2^{-\Omega(m)}$ fraction of candidate outputs, after a
constant-factor increase in output length
(\Cref{cor:projection-proofs}). This amplification concerns
certificate existence under uniform candidate outputs; the
deterministic, small-bias guarantee above is proved separately.

The small-bias guarantee supplies an explicit universal family of
candidate nonoutputs. Applying the cell-probe connection
of~\cite{KPI25,GLY26}, we obtain the following.

\begin{theorem}[Data structure consequences]\label{thm:ds-intro}
For every fixed $t\ge3$, there is an explicit Boolean query problem
with $m^{O_t(1)}$ rows and $m$ queries for which every exact
nonadaptive $t$-bit-probe representation uses
$\Omega_t(m^{2/(t-1)})$ bits of storage. In particular, three
nonadaptive probes require $\Omega(m)$ bits.
For every fixed $t\ge2$ and word size $w\ge1$, there is an explicit
query problem requiring $\Omega_{t,w}(m^{2/t})$ cells for adaptive
$t$-probe representations with $w$-bit words.
The latter bound also allows per-row relative Hamming error $1/4$.
The entries of both query problems are computable in time polynomial
in $\log m$.
\end{theorem}

The nonadaptive bound removes the logarithmic loss in the
corresponding GLY reduction, at the cost of requiring a sufficiently
small polynomial bias whose exponent is not optimized. The adaptive
word bound keeps $w$ fixed. For one-bit words, more specialized
adaptive arguments in~\cite{GLY26} achieve stronger powers than
$2/t$; we do not claim to improve those powers here.

\subsection{Relation to Previous Work}

Our range-avoidance reduction follows Guruswami, Lyu, and
Yuan~\cite[Section~6.1]{GLY26}. The change lies in the
certificate construction for arbitrary-support signed XOR systems.
The relevant deterministic thresholds are summarized below; locality
is fixed and constants may depend on it.
\begin{center}
\begin{tabular}{@{}lll@{}}
\toprule
Result & Locality & Sufficient output length\\
\midrule
Gajulapalli et al.~\cite{GGNS23} & $k\ge3$ & $n^{k-1}/\log n$\\
Guruswami--Lyu--Yuan~\cite{GLY26} & $k\ge3$ & $c^k n^{(k-1)/2}\log n$\\
\Cref{thm:avoid} & fixed $k\ge3$ & $C_k n^{(k-1)/2}$\\
\bottomrule
\end{tabular}
\end{center}

The matrix argument builds on the Kikuchi approach to semirandom
refutation and hypergraph Moore bounds~\cite{GKM22,HKM23}.
Recent proofs of the sharp hypergraph Moore bound were given by
Bandeira et al.~\cite[Theorem~1.2, version~2]{BKNPW26} and
Schmidhuber and Hastings;
we use the spectral construction of the latter work. The second
version of Bandeira et al. covers odd as well as even uniformity.
Schmidhuber and Hastings~\cite{SH26Moore} prove a sharp hypergraph
Moore bound through a normalized matrix and a lift recording the
hyperedges used an odd number of times. Their odd-uniformity argument
introduces balanced transitions and reciprocal-excess weights on four
row--label incidences. We use these constructions, with their roles
specified in the proofs below.

An even cover is a binary linear dependence among hyperedges. Its
existence alone does not give the strong-cancellation certificate in
\Cref{thm:signing}. The present argument instead uses a local capacity
condition to bound the lifted operator, and a failed trace test to
find a window where that condition fails. A separate companion paper
of Schmidhuber and Hastings~\cite{SH26Kikuchi} gives sharp
Kikuchi-hierarchy guarantees for random XOR instances. Our supports
are fixed adversarially; the signs are the variables over which the
certificate is averaged. We prove the required arbitrary-support
statement directly.

\subsection{Proof Overview}

We describe the main steps for a degree-$d$ polynomial
$P(b,x)=\sum_i a_i b_i\chi_{F_i}(x)$. The signs $b_i$ always belong
to original labels.

\paragraph{A local spectral condition.}
For even $d$, represent each clause on a slice of the Boolean cube
and normalize the resulting adjacency matrix using
$\Gamma=\diag(\text{row degrees})+\bar d I$. A lifted state
consists of a row and a parity memory of original labels. Uniform
sign averaging retains precisely the walks whose final memory is
empty. If each small window admits fractional representatives with
capacity $D$ per input variable, an orientation of the lifted edges
has incoming weight at most $(D+1)\ell$. This bounds the normalized
lifted operator by $2\sqrt{(D+1)\ell/\bar d}$.

\paragraph{Finding a failed window.}
Enumerating all logarithmic-size windows would take quasipolynomial
time. We instead inspect the averaged trace. If it is large, a
nonnegligible mass of empty-memory closed walks must visit a bad
window. Some such state has inverse-polynomial forward mass. Entries
of the signed walk matrix have nonnegative Fourier coefficients:
the coefficient of a memory is exactly the mass of paths reaching
that memory. A deterministic prefix search recovers the heavy
coefficients using a small-biased space. This finds a bad window
without constructing the full parity lift. At even degree, a failed
capacity matching gives a dense block supported on that window;
its contribution can be certified by enumerating the window's input
assignments.

\paragraph{Pairs with a common total weight.}
For odd $d$, group original clauses into equal-size bundles sharing
a center, and remove the center by one Cauchy--Schwarz inequality.
This produces pair terms with coefficients $a_i a_j b_i b_j$.
The spectral Moore construction gives port loads at most one, but
arbitrary instance weights would distort the pair polynomial.
A per-pair contention estimate allows all instances of a pair to
be scaled to one common total weight $\beta$. The exact
polynomial-to-matrix identity then survives deletion of whole pairs.

\paragraph{A failed capacity condition yields a forest.}
After pair deletion the remaining pair graph need not be a union of
cliques. We use a fractional cover condition on this graph, with
capacities on residual input variables. Its failure gives a dual
weighting. A forest-polytope argument and threshold integration turn
that weighting into a forest containing more than $D$ times as many
edges as residual variables. On a forest, the products $b_i b_j$
are independent uniform signs under uniform original signs.
The forest therefore has a small expected local certificate.
Only its pairs are deleted. The original bundle normalization and
diagonal term are retained throughout the iteration.

\paragraph{One certificate and one signing.}
The procedure stops when the remaining degree is small or the
averaged trace is small. Rational tangent-line bounds turn forest
moments and trace moments into pointwise certificates. Their Fourier
coefficient masses are polynomially bounded, so their expectation
bounds transfer to every sufficiently small-biased distribution.
We add these expectations before selecting a sign vector. This
avoids any need for independence between different forests or
between the Fourier layers of the circuit.

\subsection{Organization of the Paper}
\Cref{sec:preliminaries} gives the notation, pseudorandomness tools,
and weighted orientation inequality.
\Cref{sec:extraction} proves deterministic extraction of a bad memory
state. \Cref{sec:even} gives the even-degree certificate.
\Cref{sec:forests} proves the local-capacity alternative and its
forest certificate. \Cref{sec:odd} completes the odd-degree proof.
\Cref{sec:avoidance} assembles the signing theorem and proves the
range-avoidance algorithm. \Cref{sec:hierarchy} proves the hierarchy
tradeoff and its boundary cases.
\Cref{sec:applications} proves the one-sided tests, demi-bits
consequences, and data structure lower bounds.
\Cref{sec:moore-applications} explains what short even covers provide
directly and what additional construction is required.


% Begin preliminaries
\section{Preliminaries}
\label{sec:preliminaries}

\subsection{Local Maps, Signed Systems, and Computation}
Write $[n]=\{1,\ldots,n\}$ and identify a Boolean bit $u$ with
$(-1)^u\in\{\pm1\}$. For $F\subseteq[n]$, set
$\chi_F(x)=\prod_{v\in F}x_v$, with $\chi_\varnothing=1$.
Degree in this paper means the degree of the real multilinear Fourier
expansion on the sign cube. It is distinct from polynomial degree
over $\F_2$. A labeled degree-$d$ system is a list
$F_1,\ldots,F_M\in\binom{[n]}d$ with rational coefficients
$a_i\in[-1,1]$. The associated polynomial is
\[
 P(b,x)=\sum_{i=1}^M a_i b_i\chi_{F_i}(x),
 \qquad b\in\{\pm1\}^M,\quad x\in\{\pm1\}^n.
\]
The labels $i$ are distinct even when the sets $F_i$ coincide.
All sign memories and pseudorandom distributions refer to these labels.

A certificate for $P$ is an explicitly represented function
$\mathcal C:\{\pm1\}^M\to\Q_{\ge0}$ satisfying
$\sup_x|P(b,x)|\le\mathcal C(b)$ for every $b$.
The representation consists of rational matrices, local polynomials,
and a specified evaluation procedure. Polynomial-time construction
includes the size of this representation and the time to evaluate it.
We use the usual bit model, with rational numbers encoded by signed
integer numerators and positive denominators. The accuracy parameter
and the locality are fixed unless a parameterized statement explicitly
says otherwise. We may delete unused input variables, so polynomial
dependence on $n$ is harmless for the explicit representation.

\subsection{Small Bias and Moment Certificates}
For $\mu\in\F_2^M$, write
$\chi_\mu(b)=\prod_{i:\mu_i=1}b_i$.

\begin{definition}
A distribution $\mathcal D$ on $\{\pm1\}^M$ is
\emph{fully $\eta$-biased} if
$|\E_{b\sim\mathcal D}\chi_\mu(b)|\le\eta$ for every nonzero
$\mu\in\F_2^M$.
\end{definition}

Restrictions of a fully $\eta$-biased distribution to a subset of
coordinates remain fully $\eta$-biased. We use the following elementary
finite-field construction; small-biased spaces were developed
in~\cite{NN93,AGHP92}.

\begin{lemma}[Explicit small-biased spaces]
\label{lem:small-bias}
For every integer $h\ge1$ and rational $\eta\in(0,1)$, a fully
$\eta$-biased distribution on $h$ signs can be enumerated
deterministically in time polynomial in the input length and
$\eta^{-1}$. Its sample
space, counted with multiplicity, has size $O((h/\eta)^2)$ and all
averages of rational-valued functions are rational and can be computed
exactly by enumeration.
\end{lemma}
\begin{proof}
Choose a power of two $Q$ with
$\max\{2,h/\eta\}\le Q<2\max\{2,h/\eta\}$, and work over
$\F_Q$. Enumerate all $(a,z)\in\F_Q^2$ and output
\[
 b_j=(-1)^{\operatorname{Tr}_{\F_Q/\F_2}(az^{j-1})},
 \qquad 1\le j\le h.
\]
For a nonempty set $J\subseteq[h]$, averaging its character first
over $a$ gives the indicator that
$\sum_{j\in J}z^{j-1}=0$. This is a nonzero polynomial of degree
at most $h-1$, so the character expectation is at most
$(h-1)/Q\le\eta$. A representation of $\F_Q$ can be found by
enumerating monic degree-$\log_2Q$ binary polynomials and testing
irreducibility. Even a direct enumeration and trial-factor test is
polynomial in $Q$, which suffices here.
\end{proof}

\begin{lemma}[Bias transfer]
\label{lem:bias-transfer}
Let $H(b)=\sum_\mu c_\mu\chi_\mu(b)$ and set
$\|\widehat H\|_1=\sum_\mu|c_\mu|$. For every fully
$\eta$-biased distribution $\mathcal D$,
\begin{equation}
 \abs{\E_{\mathcal D}H-\E_{\U_M}H}
 \le\eta\|\widehat H\|_1.
 \label{eq:bias-transfer}
\end{equation}
\end{lemma}
\begin{proof}
The constant Fourier term has the same expectation under both
distributions. Bound each nonconstant character expectation by
$\eta$ and apply the triangle inequality.
\end{proof}

\begin{lemma}[Even moments]
\label{lem:moment}
Let $\xi_1,\ldots,\xi_t$ be independent uniform signs and let
$c_1,\ldots,c_t\in\R$. For every integer $r\ge1$,
\begin{equation}
 \E\left(\sum_{i=1}^t c_i\xi_i\right)^{2r}
 \le(2r-1)!!\left(\sum_i c_i^2\right)^r
 \le\left(2r\sum_i c_i^2\right)^r.
 \label{eq:even-moment}
\end{equation}
\end{lemma}
\begin{proof}
Expand the power. Only words in which each index occurs evenly
survive expectation, and their coefficient products are nonnegative.
Every such word admits a pairing of equal indices. Summing over all
$(2r-1)!!$ pairings overcounts the surviving words and yields the
first bound. The second follows from $(2r-1)!!\le(2r)^r$.
\end{proof}

For an integer $p\ge1$ and $z\ge0$, define
\begin{equation}
 R_p(z):=\frac{p-1+z}{p}\ge z^{1/p}.
 \label{eq:young-envelope}
\end{equation}
This is the tangent-line bound for the concave function $z^{1/p}$
at $z=1$. We also use, for $e>0$,
\begin{equation}
 \sqrt z\le\frac e2+\frac{z}{2e}.
 \label{eq:sqrt-envelope}
\end{equation}
Both right sides are rational for rational arguments.

\begin{lemma}[Scalar certificates]
\label{lem:scalar}
Fix rational $e\in(0,1)$. Let
$P(b,x)=\sum_{i=1}^M a_ib_i\chi_{F_i}(x)$ use at most $u$ input
variables, with $|a_i|\le1$. If
$M\ge16e^{-2}(u+1)$, put
\[
 p=2(u+1),\qquad
 \Phi(b)=\sum_{x\in\{\pm1\}^u}
          \left(\frac{P(b,x)}{eM}\right)^p,
 \qquad C(b)=eM R_p(\Phi(b)).
\]
Then $C(b)\ge\sup_x|P(b,x)|$ for every $b$. For every fully
$\eta$-biased distribution with
$\eta\le e^{2(u+1)}/(8\cdot2^u)$, one has $\E C\le eM$.
Evaluation takes time polynomial in the input and $2^u$.
\end{lemma}
\begin{proof}
Maximum-versus-sum and Eqn.~\eqref{eq:young-envelope} prove
pointwise soundness. By \Cref{lem:moment},
\[
 \E_{\U_M}\Phi
 \le2^u\left(\frac{2(u+1)}{e^2M}\right)^{u+1}
 \le2^u8^{-(u+1)}\le\frac18.
\]
Expanding the moment in the original signs gives
$\|\widehat\Phi\|_1\le2^u e^{-2(u+1)}$; reducing powers
modulo two can only decrease this upper bound. \Cref{lem:bias-transfer}
therefore gives $\E\Phi\le1/4$, and hence $\E C\le eM$.
All powers and sums are rational of polynomial bit length.
\end{proof}

\subsection{Weighted Matrices and Parity Lifts}
An undirected multigraph may have parallel edges with different
labels. Each undirected edge contributes to both symmetric adjacency
entries. Suppose its edge weights are nonnegative and its adjacency
matrix is $A_+$. Write $N$ for its number of vertices, $d(S)$ for
the weighted row degree, and
\[
 \bar d=N^{-1}\sum_Sd(S),\qquad
 \Gamma=\diag(d(S))+\bar d I.
\]
When $\bar d>0$, set
$K_+=\Gamma^{-1/2}A_+\Gamma^{-1/2}$ and
$P_+=\Gamma^{-1}A_+$. The matrix $P_+$ is substochastic, since
its row sums are $d(S)/(d(S)+\bar d)\le1$. It is similar to
$K_+$, so $\|K_+\|\le1$.

Give each edge $f$ a memory increment $h_f\in\F_2^M$.
Its signed weight is its unsigned weight multiplied by
$\chi_{h_f}(b)$. The parity lift has states $(S,\mu)$ and replaces
an edge $S\leftrightarrow S'$ by
$(S,\mu)\leftrightarrow(S',\mu+h_f)$. Memories are restricted
when a proof requires a finite truncated lift. Its edges inherit the
normalization from the base rows, rather than using its own degrees.

\begin{lemma}[Weighted orientation bound]
\label{lem:orientation}
Let $\mathcal L$ be the symmetric normalized adjacency operator of
any finite subgraph of such a lift. If its unnormalized edge weights
admit a fractional orientation with incoming weight at most $\kappa$
at every state, then
\begin{equation}
 \|\mathcal L\|\le2\sqrt{\kappa/\bar d}.
 \label{eq:orientation}
\end{equation}
A fractional orientation splits an edge into nonnegative pieces and
orients each piece independently.
\end{lemma}
\begin{proof}
This is the weighted orientation argument of
Schmidhuber--Hastings~\cite[Lemma~2.5]{SH26Moore}; we include it to
specify the normalization. Taking absolute values can only increase
the absolute quadratic form, so consider a nonnegative unit vector
$f$. An oriented piece $y\to x$ of weight $w$ contributes
$2wf(x)f(y)/\sqrt{\Gamma_x\Gamma_y}$. For any $t>0$, this is at most
\[
 w\left(t\frac{f(x)^2}{\Gamma_x}
         +t^{-1}\frac{f(y)^2}{\Gamma_y}\right).
\]
The total incoming weight is at most $\kappa$ and
$\Gamma_x\ge\bar d$. The outgoing weight at $y$ is at most its
base-row degree, which is at most $\Gamma_y$. Summing gives
$\langle f,\mathcal Lf\rangle\le t\kappa/\bar d+t^{-1}$.
Minimization over $t$ proves the bound. For $\kappa=0$ there are
no positive-weight edges and the conclusion is immediate.
\end{proof}

\begin{lemma}[Trace coefficient mass]
\label{lem:trace-mass}
Let $A_b$ be obtained from $A_+$ by multiplying each edge weight by
a sign character and a fixed real number of magnitude at most one.
Set $K_b=\Gamma^{-1/2}A_b\Gamma^{-1/2}$. For every integer $s\ge1$,
\begin{equation}
 \left\|\widehat{\tr(K_b^{2s})}\right\|_1\le N,
 \qquad
 \tr(K_b^{2s})\ge\|K_b\|^{2s}\ge0.
 \label{eq:trace-mass}
\end{equation}
Moreover, for rational data,
$\tr(K_b^{2s})=\tr((\Gamma^{-1}A_b)^{2s})$ is rational and can
be evaluated without computing matrix square roots.
\end{lemma}
\begin{proof}
Expand the trace into labeled closed walks and take absolute values
before collecting Fourier characters. Their total is at most
$\tr(K_+^{2s})\le N\|K_+\|^{2s}\le N$.
The second assertion follows from symmetry and the even power.
The rational identity follows by similarity.
\end{proof}

If a one-label move toggles a single original sign, an empty-memory
closed walk of length $2s$ has at most $s$ active labels at any
prefix. If each move toggles two labels, the corresponding bound is
$2s$: at time $j$ the memory size is at most both $2j$ and
$2(2s-j)$. These bounds concern closed walks with empty final memory;
arbitrary walk prefixes need not satisfy them.

% End preliminaries


% Begin extraction
\section{Deterministic Extraction of a Bad Memory State}
\label{sec:extraction}

The trace arguments used later in the paper identify a collection of
closed walks that pass through an obstructed state.  The underlying graph
has polynomially many vertices, but a state also records the parity of the
original labels encountered along the walk.  Enumerating all such memories
would be too expensive.  We show that a lower bound on the total weight of
the obstructed walks allows us to recover a suitable memory deterministically.
The argument uses nonnegative Fourier coefficients and applies to directed
graphs, without a reversibility assumption.

\subsection{Graphs with parity memory}

Let $G$ be a directed multigraph with vertex set $V$, where $|V|=N$.
Each edge $e$ has a nonnegative rational weight $w_e$ and an increment
$h_e\in\F_2^M$.  Loops and parallel edges are allowed, and parallel edges
retain their individual increments.  The graph is given explicitly; the
same statements apply when its edges can be enumerated in polynomial time.
Fix positive rational numbers $\gamma_v$, $v\in V$, satisfying
\begin{equation}
  \gamma_v\geq\sum_{e:\,\mathrm{tail}(e)=v}w_e.
  \label{ex:normalization}
\end{equation}
For $b\in\{\pm1\}^M$ and $\mu\in\F_2^M$, write
$\chi_\mu(b)=\prod_{i=1}^M b_i^{\mu_i}$.  Define the matrices
\begin{equation}
  A_b[v,w]=\sum_{e:\,v\to w}w_e\chi_{h_e}(b),
  \qquad
  P_b=\Gamma^{-1}A_b,
  \qquad
  \Gamma=\diag(\gamma_v:v\in V).
  \label{ex:signed-transition}
\end{equation}
We use the subscript $+$ for evaluation at the all-positive sign vector.
Thus $P_+$ is entrywise nonnegative and has row sums at most one.

The \emph{parity lift} has states $(v,\mu)\in V\times\F_2^M$.
An edge $e:v\to w$ gives, for every $\mu$, a transition
\[
  (v,\mu)\longrightarrow(w,\mu+h_e)
  \quad\text{of weight }w_e/\gamma_v.
\]
The weight of a lifted walk is the product of its transition weights.
We do not construct the parity lift explicitly.

Fix a length $T=2s$ and a predicate $\operatorname{Bad}(v,\mu)$ that can
be evaluated in polynomial time on an explicit state.  A closed
empty-memory walk starts at $(u,0)$ and returns to $(u,0)$ after $T$ steps.
Let $\Delta_{\mathrm{bad}}$ denote the total weight of such walks that
visit a bad state at some time in $\{0,\ldots,T\}$, summed over all
$u\in V$.  A walk that visits several bad states is counted only once
in this definition.  Substochasticity gives
$0\leq\Delta_{\mathrm{bad}}\leq N$.

\begin{theorem}[Deterministic memory extraction]
\label{thm:extraction}
In the setting above, suppose that a rational number $\delta>0$ satisfies
$\Delta_{\mathrm{bad}}\geq\delta$.  There is a deterministic algorithm
that finds a bad state reachable by a positive-weight lifted walk of
length at most $T$ from some empty-memory state.  Its bit complexity is
polynomial in the input description length, the predicate evaluation
time, $T$, $M$, and $1/\delta$.

More precisely, set
\begin{equation}
  \tau=\frac{\delta}{2N(T+1)}.
  \label{ex:threshold}
\end{equation}
The algorithm can enumerate a polynomial-size list containing a bad
state $(v,\mu)$ for which the total weight of $j$-step walks from
$(u,0)$ to $(v,\mu)$ is at least $\tau$, for some $u\in V$ and
$0\leq j\leq T$.
\end{theorem}

The final list may also contain states whose corresponding forward weight
lies between $\tau/2$ and $\tau$.  Every listed state is reachable by a
positive-weight walk.  This distinction is useful when an application
retains only states satisfying an additional restriction on their memories.

\subsection{Recovering nonnegative Fourier coefficients}

We first give the enumeration procedure that will be applied to entries
of powers of $P_b$.  The procedure requires exact evaluation of a function,
but does not require its Fourier expansion as input.

\begin{lemma}[Positive Fourier enumeration]
\label{lem:heavy-fourier}
Let $F:\{\pm1\}^M\to\mathbb{Q}$ have an expansion
\begin{equation}
  F(b)=\sum_{\mu\in\F_2^M}c_\mu\chi_\mu(b),
  \qquad c_\mu\geq0,
  \qquad\sum_\mu c_\mu\leq1.
  \label{ex:positive-fourier}
\end{equation}
Given an exact evaluation algorithm for $F$ and a rational
$0<\tau\leq1$, one can deterministically produce a list containing
every $\mu$ with $c_\mu\geq\tau$, and containing only $\mu$ with
$c_\mu\geq\tau/2$.  The list has size at most $2/\tau$.
The algorithm uses $\poly(M,1/\tau)$ evaluations of $F$ and polynomially
many additional bit operations in these parameters and the bit lengths
of $\tau$ and the evaluation outputs.
\end{lemma}

\begin{proof}
For a prefix $p\in\F_2^h$, define its coefficient mass by
\[
  C(p)=\sum_{\mu:\,\mu_{[h]}=p}c_\mu.
\]
We estimate $C(p)$ by assigning the first $h$ signs from a small-biased
space and fixing all remaining signs to $+1$.  In this recovery procedure
we use the separate bias parameter
\begin{equation}
  \eta_{\mathrm{rec}}=\tau/4.
  \label{ex:recovery-bias}
\end{equation}
Let $\mathcal{Q}_h$ be the explicit $\eta_{\mathrm{rec}}$-biased
distribution on $h$ signs supplied by \Cref{lem:small-bias}.  Set
\begin{equation}
  \widehat C(p)=
  \E_{z\sim\mathcal{Q}_h}
    \bigl[\chi_p(z)F(z,1,\ldots,1)\bigr].
  \label{ex:prefix-oracle}
\end{equation}
This expectation is computed exactly by enumerating the sample space.
Indeed, substituting Eqn.~\eqref{ex:positive-fourier} into
Eqn.~\eqref{ex:prefix-oracle} gives
\[
  \widehat C(p)=
  \sum_\mu c_\mu
    \E_{z\sim\mathcal{Q}_h}\chi_{p+\mu_{[h]}}(z).
\]
The summands with $\mu_{[h]}=p$ contribute exactly $C(p)$.
Every other character has expectation of absolute value at most
$\eta_{\mathrm{rec}}$.  Consequently,
\begin{equation}
  \abs{\widehat C(p)-C(p)}
    \leq\eta_{\mathrm{rec}}\sum_\mu c_\mu
    \leq\tau/4.
  \label{ex:prefix-error}
\end{equation}
This is also an instance of \Cref{lem:bias-transfer}.

Starting with the empty prefix, extend every retained prefix by both
possible bits at the next coordinate.  Retain an extension $p$ if
\begin{equation}
  \widehat C(p)\geq3\tau/4.
  \label{ex:prefix-test}
\end{equation}
If $c_\mu\geq\tau$, then every prefix of $\mu$ has mass at least
$\tau$, so Eqn.~\eqref{ex:prefix-error} ensures that all these prefixes
are retained.  Conversely, any retained prefix has mass at least
$\tau/2$.  Prefix masses at a fixed depth concern disjoint sets of
coefficients and sum to at most one.  Thus at most $2/\tau$ prefixes
are retained at each depth.  At depth $M$, a prefix specifies an
individual coefficient, proving both guarantees on the output list.

There are $M$ depths and at most $4/\tau$ oracle calls at each depth.
Each oracle call enumerates a space of size
$\poly(M,1/\eta_{\mathrm{rec}})=\poly(M,1/\tau)$.
All comparison thresholds are rational, and the expectations are exact
rational averages.  The claimed complexity follows.
\end{proof}

Fixing the suffix coordinates to $+1$ in
Eqn.~\eqref{ex:prefix-oracle} is essential to this proof: it makes the
oracle estimate the sum of all nonnegative coefficients extending a
prefix.  No individual coefficient needs to be identified before its
entire index has been specified.

\subsection{Forward and suffix masses}

For $u,v\in V$ and $0\leq j\leq T$, consider the efficiently
evaluable function
\[
  F_{u,v,j}(b)=(P_b^j)[u,v].
\]
Expanding matrix multiplication and reducing characters using
$\chi_\mu\chi_\nu=\chi_{\mu+\nu}$ yields
\begin{equation}
  F_{u,v,j}(b)=
  \sum_{\mu\in\F_2^M}c_{u,v,j}(\mu)\chi_\mu(b).
  \label{ex:path-fourier}
\end{equation}
Here $c_{u,v,j}(\mu)$ is exactly the total weight of $j$-step lifted
walks from $(u,0)$ to $(v,\mu)$.  In particular, these Fourier
coefficients are nonnegative.  Moreover,
\begin{equation}
  \sum_\mu c_{u,v,j}(\mu)=(P_+^j)[u,v],
  \qquad
  \sum_{v,\mu}c_{u,v,j}(\mu)\leq1.
  \label{ex:path-mass}
\end{equation}
Thus \Cref{lem:heavy-fourier} applies to every $F_{u,v,j}$.

The following argument explains why at least one recoverable coefficient
describes a bad state.  Translation invariance of the memory coordinate
implies that the total weight of $r$-step walks from $(v,\mu)$ to
$(u,0)$ equals $c_{v,u,r}(\mu)$: in both cases the required memory
increment is $\mu$.  This identity does not reverse the direction of
any edge.  It follows that the total weight of closed empty-memory
$T$-step walks starting at $(u,0)$ and passing through $(v,\mu)$
at time $j$ is
\begin{equation}
  c_{u,v,j}(\mu)c_{v,u,T-j}(\mu).
  \label{ex:prefix-suffix}
\end{equation}

For a fixed time $j$, call a visit \emph{light} relative to its starting
vertex $u$ if $c_{u,v,j}(\mu)<\tau$.  Summing
Eqn.~\eqref{ex:prefix-suffix} over all light visits gives
\begin{align}
 &\sum_{u\in V}
   \sum_{\substack{v\in V,\,\mu\in\F_2^M\\
                    c_{u,v,j}(\mu)<\tau}}
       c_{u,v,j}(\mu)c_{v,u,T-j}(\mu)
       \notag\\
 &\hspace{20mm}\leq
       \tau\sum_{v\in V}\sum_{u\in V,\,\mu\in\F_2^M}
          c_{v,u,T-j}(\mu)
       \leq N\tau.
  \label{ex:light-mass}
\end{align}
The last inequality is the row-sum bound in
Eqn.~\eqref{ex:path-mass}, applied with $v$ as the starting vertex
of the suffix.  In particular, no estimate on a column sum or on the
ratios $\gamma_u/\gamma_v$ is needed.

\begin{proof}[Proof of \Cref{thm:extraction}]
Consider the bad closed walks for which every bad visit is light.
Each such walk has at least one light bad visit.  Summing
Eqn.~\eqref{ex:light-mass} over the $T+1$ possible visit times
therefore bounds their total weight by
\[
  N(T+1)\tau=\delta/2.
\]
Since $\Delta_{\mathrm{bad}}\geq\delta$, there exists a bad
visit $(v,\mu)$ at a time $j$ on a closed walk starting at $(u,0)$
such that $c_{u,v,j}(\mu)\geq\tau$.

Apply \Cref{lem:heavy-fourier}, with the threshold in
Eqn.~\eqref{ex:threshold}, to each of the $N^2(T+1)$ functions
$F_{u,v,j}$.  Retain the indices $u,j$ along with each recovered
memory $\mu$ and endpoint $v$, and evaluate
$\operatorname{Bad}(v,\mu)$ on the resulting list.  The bad visit
just established appears in the list, so at least one test succeeds.
Every listed coefficient is at least $\tau/2>0$ and therefore
corresponds to a reachable state.

For completeness, all evaluations and comparisons admit polynomial
bit complexity.  Constructing $P_b$ uses the explicitly given rational
edge weights and normalizers.  A common denominator for its entries
has bit length polynomial in the graph description.  Powers up to
$T$ have denominators dividing the $T$-th power of such a denominator,
and their numerators also have polynomial bit length.  They can be
computed by exact rational matrix multiplication.  The prefix oracles
use polynomially many evaluations and exact finite averages; their
bit lengths remain polynomial in the stated parameters.  Finally,
$1/\tau=2N(T+1)/\delta$, so the enumeration complexity in
\Cref{lem:heavy-fourier} has the required dependence.
\end{proof}

\subsection{Truncation in the applications}

The preceding theorem allows arbitrary increments.  In our applications,
each increment toggles either one original label or a pair of distinct
original labels.  The following observation permits recovery to be
followed by a memory-size filter.

Suppose every increment has Hamming weight at most $a$, and consider a
closed empty-memory walk of length $T=2s$.  If its memory after $j$
steps is $\mu_j$, then both its prefix and its suffix bound that
memory:
\begin{equation}
  \abs{\supp(\mu_j)}
     \leq\min\{aj,a(T-j)\}
     \leq as.
  \label{ex:memory-filter}
\end{equation}
Consequently, the heavy bad visit supplied by the proof of
\Cref{thm:extraction} survives a filter discarding all memories
larger than $as$.  The fact that other recovered states need not lie
on a closed walk causes no difficulty.

For the paired construction, $a=2$.  A state consists of an
$\ell$-subset $S$ of the input variables and a memory $\mu$ of
original labels.  If each label has a residual set $R_i$ of size at
most $d$, its associated window is
\[
  U(S,\mu)=S\cup\bigcup_{i\in\supp(\mu)}R_i.
\]
After the filter in Eqn.~\eqref{ex:memory-filter}, this window has
size at most $\ell+2ds$.  Hence a polynomial-time feasibility test
on an explicit window can serve as the bad-state predicate, even
though enumerating all windows of this size would not generally
take polynomial time.  When $N$ is polynomial in the input length,
$s=O(\log N)$, and $\delta^{-1}$ is polynomial, the entire
extraction procedure has polynomial bit complexity.

% End extraction


% Begin even
\section{Certificates for Even-Degree XOR}\label{sec:even}

We first construct certificates for even-degree polynomials. The proof
combines the normalized parity lift of the spectral Moore argument
with a decomposition into small dense windows. The orientation argument
is adapted from Schmidhuber and Hastings~\cite[Lemmas~2.5--2.7]{SH26Moore};
the decomposition removes its local sparsity hypothesis. All random
characters below are characters of the original clause labels, including
when several labels have identical scopes.

\begin{theorem}\label{thm:even-weight}
Fix an even integer $d\geq2$ and a rational $0<\epsilon<1$.
Let $F_1,\ldots,F_M$ be labeled $d$-subsets of $[n]$, where $n\geq d$,
and let $a_1,\ldots,a_M$ be rational numbers with $\abs{a_i}\leq1$.
Write $W=\sum_i\abs{a_i}$. There is a constant $C_{d,\epsilon}$ such that,
if $W\geq C_{d,\epsilon}n^{d/2}$, a deterministic algorithm constructs
a nonnegative rational-valued function $\operatorname{Cert}$ satisfying
\[
 \sup_{x\in\{\pm1\}^n}
 \abs{\sum_{i=1}^M a_i b_i\chi_{F_i}(x)}
 \leq\operatorname{Cert}(b)
 \qquad\text{for every }b\in\{\pm1\}^M.
\]
For an explicitly specified $\eta=n^{-O_{d,\epsilon}(1)}>0$, it also satisfies
\begin{equation}\label{even:main-expectation}
 \E_{b\sim\mathcal D}\operatorname{Cert}(b)
 \leq \frac{41}{128}\epsilon W
\end{equation}
for every $\eta$-biased distribution $\mathcal D$ on the original $M$ signs.
Construction and evaluation have polynomial bit complexity for fixed
$d,\epsilon$.
\end{theorem}

The decomposition used to construct $\operatorname{Cert}$ may depend on a
fixed explicit biased space. The distribution $\mathcal D$ in
Eqn.~\eqref{even:main-expectation} is arbitrary and need not be that space.
This distinction will allow different Fourier layers to share one final
sign vector.

\subsection{Normalized matrices and local allocations}

Set $\ell=d/2$, $N=\binom n\ell$, and $c=\binom d\ell$.
For each label $i$, let $A_i$ be the adjacency matrix on
$\Omega=\binom{[n]}\ell$ that joins $S$ to $S\triangle F_i$ whenever
$\abs{S\cap F_i}=\ell$. Each undirected edge is included once, so
$A_i$ has exactly $c$ nonzero entries. Put $w_i=\abs{a_i}$ and discard
zero-weight labels from the matrices, while retaining their positions in
the original sign vector.

For a nonempty set $R$ of positive-weight labels, write
\[
 W_R=\sum_{i\in R}w_i,\qquad
 A_R^{\mathrm{abs}}=\sum_{i\in R}w_i A_i,\qquad
 A_{R,b}=\sum_{i\in R}a_i b_i A_i.
\]
Let $d_R(S)$ be the row sum of $A_R^{\mathrm{abs}}$, and define
\begin{equation}\label{even:normalization}
 \overline d_R=\frac{cW_R}{N},\qquad
 \Gamma_R=\diag(d_R(S))+\overline d_R I,\qquad
 K_{R,b}=\Gamma_R^{-1/2}A_{R,b}\Gamma_R^{-1/2}.
\end{equation}
For $z(x)_S=\chi_S(x)$, direct counting gives
\begin{equation}\label{even:comparison}
 z(x)^T A_{R,b}z(x)=c\sum_{i\in R}a_i b_i\chi_{F_i}(x),
 \qquad z(x)^T\Gamma_Rz(x)=2cW_R.
\end{equation}
Consequently, $2W_R\norm{K_{R,b}}$ is a pointwise upper bound for the
absolute value of the residual polynomial.

Fix a capacity $D>0$ and a window $V\subseteq[n]$. An allocation in $V$
consists of nonnegative numbers $\lambda_{i,v}$, for labels $i\in R$
with $F_i\subseteq V$ and vertices $v\in F_i$, satisfying
\begin{equation}\label{even:allocation}
 \sum_{v\in F_i}\lambda_{i,v}=w_i,
 \qquad
 \sum_{i:F_i\subseteq V}\lambda_{i,v}\leq D
 \quad(v\in V).
\end{equation}
Feasibility is a rational maximum-flow problem: labels have demands $w_i$,
vertices have capacities $D$, and a label can send flow only to its scope.
The flow condition is
$\sum_{i\in J}w_i\leq D\abs{\bigcup_{i\in J}F_i}$ for every subfamily
$J$ of the labels in $V$. Thus a failed flow supplies a subwindow
$U\subseteq V$ such that
\begin{equation}\label{even:dense-window}
 \sum_{i\in R:F_i\subseteq U}w_i>D\abs U.
\end{equation}
Indeed, take $U$ to be the union of a violating subfamily. Both the flow
test and the extraction of $U$ take polynomial time.

\begin{lemma}\label{even:good-trace}
Fix an integer $s\geq1$. In the parity lift truncated to memories of
size at most $s$, retain only states whose windows admit
Eqn.~\eqref{even:allocation}. The total weight of its closed
empty-memory walks of length $2s$, summed over starting rows, is at most
\begin{equation}\label{even:good-mass}
 N\left(\frac{4(D+1)\ell}{\overline d_R}\right)^s.
\end{equation}
The total weight before removing any states equals
$\E_{b\text{ uniform}}\tr(K_{R,b}^{2s})$.
\end{lemma}

\begin{proof}
A lifted state is $(S,P)$, where $S\in\Omega$ and $P\subseteq R$ has
$\abs P\leq s$. A move of label $i$ changes it to
$(S\triangle F_i,P\triangle\{i\})$ and has unnormalized weight $w_i$.
Its normalized weight is $w_i/\sqrt{\Gamma_R(S,S)\Gamma_R(T,T)}$.
The window of $(S,P)$ is
\[
 V(S,P)=S\cup\bigcup_{i\in P}F_i.
\]
Along an edge the two windows are nested: inserting $i$ into memory
replaces the window by its union with $F_i$.

Fix one feasible allocation for each good window. For an edge whose
endpoint windows agree, split its weight into pieces $\lambda_{i,v}$
and orient each piece toward the endpoint whose row contains $v$.
The row toggle ensures that exactly one endpoint contains $v$.
At a fixed state $(S,P)$, these pieces have total incoming weight at most
$D\ell$: at each $v\in S$ they are a subset of that window's allocation,
and at most one incident edge has any given original label.

Orient an edge with unequal windows toward the larger window.
For such an edge entering $(S,P)$, choose a vertex $v$ newly present
in its window. The toggled label $i$ was inserted into memory.
Since $v$ belongs to $F_i$ but not to the preceding row, the row toggle
places it in $S$. Moreover, $i$ is the unique label in $P$ containing
$v$, because any other such label was already present at the other
endpoint. A witness $v$ therefore determines the incoming edge.
There are at most $\ell$ choices of witness, and each edge weighs at most
one. The total incoming weight is at most $(D+1)\ell$.
The outgoing weight at a state with row $S$ is at most $d_R(S)$.
By \Cref{lem:orientation}, the induced good-state operator has norm
at most $2\sqrt{(D+1)\ell/\overline d_R}$. Its return mass from the
$N$ empty-memory states is bounded by Eqn.~\eqref{even:good-mass}.

Finally, expand $\tr(K_{R,b}^{2s})$ into labeled closed walks.
Uniform signs kill every word in which some label occurs oddly.
In each surviving word, the product of the $a_i$ equals the product of
the $w_i$, since all multiplicities are even. Such a word lifts uniquely
to an empty-memory closed walk. At time $j$ its memory has size at most
$\min\{j,2s-j\}\leq s$, so truncation loses none of these words.
This proves the trace identity, without any assumption about even covers.
\end{proof}

For later use, the absolute Fourier coefficient mass of
$\tr(K_{R,b}^{2s})$ is at most $N$.
Indeed, the sum of absolute walk contributions is
$\tr((K_R^{\mathrm{abs}})^{2s})$, where
$K_R^{\mathrm{abs}}=\Gamma_R^{-1/2}A_R^{\mathrm{abs}}\Gamma_R^{-1/2}$.
This matrix is similar to the substochastic matrix
$\Gamma_R^{-1}A_R^{\mathrm{abs}}$, hence has norm at most one.
In particular, \Cref{lem:bias-transfer} gives
\begin{equation}\label{even:trace-transfer}
 \abs{\E_{b\sim\mathcal D}\tr(K_{R,b}^{2s})
       -\E_{b\text{ uniform}}\tr(K_{R,b}^{2s})}
 \leq\eta N
\end{equation}
for every $\eta$-biased $\mathcal D$.

\subsection{A deterministic decomposition}

Choose the following parameters:
\begin{align}
 s&=\lceil\log_2(4N)\rceil,&
 L&=\min\{n,\ell+ds\},&
 \zeta&=\epsilon/32,\notag\\
 D&=(64/\epsilon)^2,&
 \beta&=\zeta^{2s},&
 \eta&=\min\left\{\frac{\beta}{8N},
                    (\epsilon/64)^{2L}\right\}.
 \label{even:parameters}
\end{align}
It is enough to impose the explicit density condition
\begin{equation}\label{even:density}
 W\geq\frac{64(D+1)\ell}{\epsilon\zeta^2}\frac Nc.
\end{equation}
This follows from $W\geq C_{d,\epsilon}n^{d/2}$ after choosing the
constant sufficiently large. Fix one explicit $\eta$-biased space
$\mathcal Q$ on the original $M$ coordinates, as in
\Cref{lem:small-bias}.

Start with all positive-weight labels in $R$ and an empty list of blocks.
Perform the following steps until a remainder is recorded.
\begin{enumerate}
\item If $W_R\leq\epsilon W/8$, record $R$ as a small remainder and stop.
\item Otherwise form the matrices in Eqn.~\eqref{even:normalization}
and compute the exact rational average
\[
 T_{\mathcal Q}=
 \E_{b\sim\mathcal Q}\tr\bigl((\Gamma_R^{-1}A_{R,b})^{2s}\bigr).
\]
If $T_{\mathcal Q}\leq\beta$, record $R$ as a spectral remainder and stop.
\item If $T_{\mathcal Q}>\beta$, find a window $U$ satisfying
Eqn.~\eqref{even:dense-window}, with $\abs U\leq L$.
Record all labels of $R$ supported on $U$ as a block, remove them from
$R$, and repeat.
\end{enumerate}

We justify the third step. The first test has failed, and
Eqn.~\eqref{even:density} therefore implies
\[
 \overline d_R=\frac{cW_R}{N}
 >\frac{8(D+1)\ell}{\zeta^2}.
\]
By Eqn.~\eqref{even:trace-transfer}, the uniform-sign trace is greater
than $7\beta/8$. On the other hand, \Cref{even:good-trace} bounds the
mass of empty-memory closed walks staying in good states by
\[
 N\left(\frac{\zeta^2}{2}\right)^s\leq\frac\beta4.
\]
The mass visiting a bad state is consequently greater than $5\beta/8$.
Apply \Cref{thm:extraction}, with walk length $2s$ and
$\delta=\beta/2$, to the nonnegative labeled transition matrix
$\Gamma_R^{-1}A_R^{\mathrm{abs}}$. Its memory increment for label $i$
is the $i$th original coordinate. Closed transition weights agree with
the symmetric normalized weights by telescoping the diagonal factors.
The bad predicate first rejects memories larger than $s$ and then tests
allocation feasibility in their windows. Every relevant closed walk has
memory at most $s$, so this filter loses none of the promised bad mass.
The extracted window has size at most $L$, and its failed flow supplies
the required subwindow $U$.

Each iteration removes at least one positive-weight label. There are
therefore at most $M$ iterations, and the recorded blocks are disjoint
in original labels. Their windows may overlap.

\subsection{Pointwise certificates and their expectations}

Let $I_j$ be a recorded block with window $U_j$, and put
$W_j=\sum_{i\in I_j}w_i$. Define its exact local certificate by
\[
 C_j(b)=\max_{x\in\{\pm1\}^{U_j}}
       \abs{\sum_{i\in I_j}a_i b_i\chi_{F_i}(x)}.
\]
A small remainder contributes $C_R(b)=W_R$.
For a spectral remainder, put
\[
 T_R(b)=\tr\bigl((\Gamma_R^{-1}A_{R,b})^{2s}\bigr),
 \qquad \alpha_R(b)=T_R(b)^{1/(2s)}.
\]
The trace is nonnegative because the matrix is similar to a real
symmetric matrix. Compute a rational $\widehat\alpha_R(b)$ with
\[
 \alpha_R(b)\leq\widehat\alpha_R(b)
 \leq\alpha_R(b)+\epsilon/256,
\]
and set $C_R(b)=2W_R\widehat\alpha_R(b)$.
Finally, define
\begin{equation}\label{even:certificate}
 \operatorname{Cert}(b)=\sum_j C_j(b)+C_R(b).
\end{equation}
The local certificates are exact. For a spectral remainder,
$\norm{K_{R,b}}\leq\alpha_R(b)$, so
Eqn.~\eqref{even:comparison} proves its pointwise soundness.
The triangle inequality now proves pointwise soundness of
Eqn.~\eqref{even:certificate} for every sign vector.

Fix an arbitrary final $\eta$-biased distribution $\mathcal D$.
For a recorded block, set $t=\abs{U_j}\geq1$ and
$p_b(x)=\sum_{i\in I_j}a_i b_i\chi_{F_i}(x)$.
The Rademacher moment bound (\Cref{lem:moment}) and
$\sum_{i\in I_j}a_i^2\leq W_j$ give
\[
 \E_{b\text{ uniform}}p_b(x)^{2t}
 \leq(2t-1)!!\,W_j^t\leq(2tW_j)^t.
\]
As a polynomial in the original signs, $(p_b(x)/W_j)^{2t}$ has
Fourier coefficient mass at most one: expanding before reducing
$b_i^2=1$ bounds that mass by
$(\sum_{i\in I_j}\abs{a_i}/W_j)^{2t}=1$.
Thus \Cref{lem:bias-transfer} implies
\[
 \E_{b\sim\mathcal D}p_b(x)^{2t}
 \leq(2tW_j)^t+\eta W_j^{2t}.
\]
Use $C_j(b)^{2t}\leq\sum_x p_b(x)^{2t}$, Jensen's inequality,
and subadditivity of the $1/(2t)$ power to obtain
\begin{align}
 \E_{b\sim\mathcal D}C_j(b)
 &\leq 2\sqrt{tW_j}+\sqrt2\,\eta^{1/(2t)}W_j\notag\\
 &\leq\left(\frac{2}{\sqrt D}
              +\frac{\sqrt2\epsilon}{64}\right)W_j
 \leq\frac{\epsilon}{16}W_j.
 \label{even:block-expectation}
\end{align}
Here $W_j>Dt$ by extraction, and $t\leq L$ with
Eqn.~\eqref{even:parameters} gives
$\eta^{1/(2t)}\leq\epsilon/64$.

A small remainder contributes at most $\epsilon W/8$.
At a spectral stop, compare the final distribution to uniform and
then to $\mathcal Q$ using Eqn.~\eqref{even:trace-transfer}:
\[
 \E_{b\sim\mathcal D}T_R(b)
 \leq\beta+2\eta N\leq\frac54\beta.
\]
Jensen's inequality gives
$\E\alpha_R(b)\leq((5/4)\beta)^{1/(2s)}\leq2\zeta$, and hence
\begin{equation}\label{even:remainder-expectation}
 \E_{b\sim\mathcal D}C_R(b)
 \leq4\zeta W_R+\frac{\epsilon}{128}W_R
 \leq\left(\frac18+\frac1{128}\right)\epsilon W.
\end{equation}
The blocks are disjoint in labels, so $\sum_jW_j\leq W$.
Charging both possible remainder bounds, although only one occurs,
Eqns.~\eqref{even:block-expectation} and
\eqref{even:remainder-expectation} imply
\[
 \E_{b\sim\mathcal D}\operatorname{Cert}(b)
 \leq\left(\frac1{16}+\frac18+\frac18+\frac1{128}\right)
       \epsilon W
 =\frac{41}{128}\epsilon W.
\]
No independence between blocks is used in summing these expectations.

\subsection{Bit complexity}

For fixed $d,\epsilon$, we have $N=n^{O_d(1)}$, $s,L=O_d(\log n)$,
and $\beta^{-1},\eta^{-1}=n^{O_{d,\epsilon}(1)}$.
The explicit biased space has size polynomial in $M$ and $1/\eta$.
The decomposition uses at most $M$ iterations, and
\Cref{thm:extraction} makes each window search polynomial in the graph
input size and $1/\beta$. Every recorded local maximum can be evaluated
by enumerating $2^{\abs{U_j}}\leq2^L$ assignments.

All matrices, capacities, and trace averages are rational.
The trace is computed using $\Gamma_R^{-1}A_{R,b}$, so no matrix square
roots are needed. A common denominator has polynomial bit length, and
powers through $2s$ preserve a polynomial bit bound.
Moreover, $T_R(b)\leq N$ by absolute domination, so
$\alpha_R(b)\leq N^{1/(2s)}<2$.
Bisection on $[0,2]$, with comparisons of rational $2s$th powers,
therefore computes the required upper approximation with polynomial bit
complexity. This completes the proof of \Cref{thm:even-weight}.

% End even


% Begin forests
\section{Dense Forests from Local Capacity Violations}
\label{sec:forests}

The odd-degree construction will delete pairs of original labels while
retaining the original sign variables.  A forest of such pairs is useful
because its edge products are independent under uniform original signs.
The first lemma obtains a dense forest from a failed fractional capacity
condition.  The second turns this forest into a rational certificate whose
expectation is small under every sufficiently small-biased distribution on the
original signs.

\begin{lemma}[A forest from a capacity violation]
\label{lem:forest-alternative}
Let $G=(V,E)$ be a finite simple graph, let each $i\in V$ have a nonempty
residual set $R_i$ in a finite universe $U$, and let $D\geq 0$ be rational.
Consider the system
\begin{equation}
\begin{aligned}
 \lambda_{iv}&\geq 0 &&(i\in V,\ v\in R_i),\\
 t_i&=\sum_{v\in R_i}\lambda_{iv} &&(i\in V),\\
 t_i+t_j&\geq 1 &&(ij\in E),\\
 \sum_{i:v\in R_i}\lambda_{iv}&\leq D &&(v\in U).
\end{aligned}
\label{forest:eq:capacity}
\end{equation}
If this system is infeasible, one can find a forest $T\subseteq E$ such
that, writing $V(T)$ for its endpoints and
$U_T=\bigcup_{i\in V(T)}R_i$,
\begin{equation}
 |T|>D|U_T|.
 \label{forest:eq:density}
\end{equation}
The forest can be found deterministically in polynomial bit complexity.
Moreover, adding the conditions $t_i\leq 1$ to
\eqref{forest:eq:capacity} does not change its feasibility.
\end{lemma}

\begin{proof}
First suppose a feasible allocation has $t_i>1$.  Dividing all
$\lambda_{iv}$ at this vertex by $t_i$ reduces its total allocation to
one and decreases every capacity load.  Each incident edge remains
covered because its endpoint $i$ alone now contributes one.  Applying
this operation at every such vertex proves the last assertion.

Farkas' lemma gives nonnegative rational edge weights $y_e$ and vertex
prices $p_v$, for $v\in U$, such that
\begin{equation}
 d_i:=\sum_{e\ni i}y_e\leq p_v\quad(v\in R_i),
 \qquad
 \sum_{e\in E}y_e>D\sum_{v\in U}p_v.
 \label{forest:eq:dual}
\end{equation}
Indeed, the coefficient inequality for the nonnegative variable
$\lambda_{iv}$ is precisely $d_i\leq p_v$.  Since the residual sets
are nonempty, any witness has $\sum_vp_v>0$.  We may therefore normalize
it to $\sum_vp_v=1$.  An explicit way to find a witness is to maximize
$\sum_ey_e$ subject to this normalization, nonnegativity, and the
inequalities $d_i\leq p_v$.  This is a bounded rational linear program:
each $d_i\leq 1$, and hence each $y_e\leq 1$.  Its optimum is greater
than $D$ whenever \eqref{forest:eq:capacity} is infeasible.

For the threshold argument, replace each price by its smallest allowed
value,
\begin{equation}
 p_v=\max_{i:v\in R_i}d_i,
 \label{forest:eq:minimal-prices}
\end{equation}
where an empty maximum is zero.  This preserves the strict inequality
in \eqref{forest:eq:dual}.  If desired, rescale both $y$ and $p$ once
more so that the price sum is one; this also preserves
\eqref{forest:eq:minimal-prices}.

For $e=ij$, put $w_e=\min(d_i,d_j)$ and define
\begin{equation}
 z_e=
 \begin{cases}
 y_e/w_e,&y_e>0,\\
 0,&y_e=0.
 \end{cases}
 \label{forest:eq:fractional-forest}
\end{equation}
The denominator is positive whenever it is used.  Order the vertices
by increasing degree $d_i$, breaking ties by a fixed total order, and
orient each positive-$y$ edge towards its larger endpoint.  At each
vertex of positive degree, the outgoing $z$-weight is at most one:
\[
 \sum_{e\text{ directed out of }i}z_e
 =\frac{1}{d_i}\sum_{e\text{ directed out of }i}y_e\leq 1.
\]
For every nonempty $S\subseteq V$, its largest vertex has no outgoing
edge inside $G[S]$.  Summing the outgoing bounds over the other
vertices gives
\begin{equation}
 \sum_{e\in E(G[S])}z_e\leq |S|-1.
 \label{forest:eq:polytope}
\end{equation}
These are the graphic forest-polytope inequalities.  More explicitly,
for any edge set $Q\subseteq E$, apply
\eqref{forest:eq:polytope} to the vertex sets of its nontrivial
connected components to obtain
\begin{equation}
 \sum_{e\in Q}z_e\leq\rank(Q).
 \label{forest:eq:rank-inequality}
\end{equation}
Here $\rank(Q)$ denotes the number of edges in any spanning forest of
$(V,Q)$.

We record directly the weighted consequence of these inequalities.
For nonnegative edge weights $w$, write
$E_t=\{e:w_e\geq t\}$.  Kruskal's maximum-spanning-forest algorithm
produces a forest $T_*$ whose edges in $E_t$ form a spanning forest of
$(V,E_t)$, apart from immaterial threshold endpoints.  Therefore
\begin{equation}
 \sum_e w_ez_e
 =\int_0^\infty\sum_{e\in E_t}z_e\,dt
 \leq\int_0^\infty\rank(E_t)\,dt
 =\sum_{e\in T_*}w_e.
 \label{forest:eq:weighted-domination}
\end{equation}
This also supplies a direct justification of the needed forest-polytope
property, without appealing to an integrality theorem for a linear
program.  Applied to the weights in
\eqref{forest:eq:fractional-forest}, it gives
$\sum_{e\in T_*}w_e\geq\sum_ey_e$.

For $t>0$, let $S_t=\{i:d_i\geq t\}$ and
$U_t=\bigcup_{i\in S_t}R_i$.  The condition $w_{ij}\geq t$ is exactly
$i,j\in S_t$, while \eqref{forest:eq:minimal-prices} gives the second
of the following identities:
\begin{equation}
 \sum_{e\in T_*}w_e
 =\int_0^\infty\rank(G[S_t])\,dt,
 \qquad
 \sum_vp_v=\int_0^\infty|U_t|\,dt.
 \label{forest:eq:threshold-integrals}
\end{equation}
Combining \eqref{forest:eq:dual},
\eqref{forest:eq:weighted-domination}, and
\eqref{forest:eq:threshold-integrals}, we obtain
\[
 \int_0^\infty\bigl(\rank(G[S_t])-D|U_t|\bigr)\,dt>0.
\]
Consequently some positive degree threshold satisfies
$\rank(G[S_t])>D|U_t|$.  Take any spanning forest of $G[S_t]$ and
discard its isolated vertices.  Its endpoint residual union is
contained in $U_t$, so this forest satisfies
\eqref{forest:eq:density}.

All linear programs above have polynomially many variables and
constraints and rational data of polynomial bit length.  Standard
rational linear programming returns an optimal solution of polynomial
bit length.  There are at most $|V|$ distinct positive degree
thresholds, and the graph ranks and spanning forests at all of them
can be computed in polynomial time.  Comparisons at these rational
thresholds are exact.  This proves the asserted constructivity.
\end{proof}

The preceding implication is one-way: a feasible allocation need not
exclude dense forests.  For example, when every residual is the same
singleton, a star has a feasible allocation of capacity one by placing
all allocation at its center, however many leaves the star has.  We
will only use infeasibility to produce a forest.

\begin{lemma}[A certificate for a dense forest]
\label{lem:forest-certificate}
Let $T$ be a forest whose vertices are distinct original labels in
$[M]$, with residual sets $R_i$, and let
$U_T=\bigcup_{i\in V(T)}R_i$ have size $1\leq u\leq L$.
Let $a_i\in\mathbb Q\cap[-1,1]$, and put $m_T=|T|$ and
\[
 Q_T(b,x)=\sum_{ij\in T}a_ia_jb_ib_j
             \chi_{R_i\triangle R_j}(x).
\]
Suppose $\sigma\in\mathbb Q\cap(0,1)$,
$D\geq32\sigma^{-2}$, and $m_T>Du$.  If
\begin{equation}
 \eta\leq\frac{\sigma^{2(L+1)}}{8\cdot2^L},
 \label{forest:eq:bias}
\end{equation}
then there is a nonnegative rational certificate $C_T(b)$ such that
\begin{equation}
 \sup_{x\in\{\pm1\}^{U_T}}|Q_T(b,x)|\leq C_T(b)
 \quad\text{for every }b,
 \qquad
 \E_{b\sim\mathcal Q}C_T(b)\leq\sigma m_T
 \label{forest:eq:certificate-guarantees}
\end{equation}
for every fully $\eta$-biased distribution $\mathcal Q$ on the
original signs $b_1,\ldots,b_M$.  The certificate is exactly evaluable
in time polynomial in the input length and $2^u$.
\end{lemma}

\begin{proof}
Set $p=2(u+1)$ and define
\begin{equation}
 \Phi_T(b)=\sum_{x\in\{\pm1\}^{U_T}}
       \left(\frac{Q_T(b,x)}{\sigma m_T}\right)^p,
 \qquad
 C_T(b)=\sigma m_T R_p(\Phi_T(b)),
 \label{forest:eq:potential}
\end{equation}
where $R_p$ is the rational envelope from
\eqref{eq:young-envelope}.  The exponent is even, so every summand is
nonnegative.  Its largest summand gives
\[
 \sup_x|Q_T(b,x)|
 \leq\sigma m_T\Phi_T(b)^{1/p}
 \leq\sigma m_T R_p(\Phi_T(b))=C_T(b).
\]

Under uniform independent original signs, the products
$\{b_ib_j:ij\in T\}$ are mutually independent uniform signs.  To
see this, prescribe an arbitrary sign on every edge and choose one
root sign in each connected component.  Every other vertex sign is
then uniquely recovered along the component's tree paths.  Thus
every edge-sign assignment has exactly two preimages per connected
component, with all unused original signs remaining free.  Repeated
residual sets or repeated characters cause no change to this argument:
the vertices of the forest are original labels.

For a fixed $x$, the independent-sign moment bound in
Lemma~\ref{lem:moment}, with $t=u+1$ and coefficients of magnitude
at most one, gives
\[
 \E_{\mathrm{unif}}Q_T(b,x)^{2(u+1)}
 \leq\bigl(2(u+1)m_T\bigr)^{u+1}.
\]
Consequently
\begin{align}
 \E_{\mathrm{unif}}\Phi_T
 &\leq 2^u
       \left(\frac{2(u+1)}{\sigma^2m_T}\right)^{u+1}
 \nonumber\\
 &\leq 2^u
       \left(\frac{u+1}{16u}\right)^{u+1}
 \leq 2^u8^{-(u+1)}
 \leq\frac18.
 \label{forest:eq:uniform-moment}
\end{align}
We have used $m_T>Du$, $D\geq32\sigma^{-2}$, and $u\geq1$.

It remains to transfer this estimate on the original sign space.
For each fixed $x$, the sum of the absolute Fourier coefficients of
$Q_T(b,x)$ as a polynomial in $b$ is at most
$\sum_{ij\in T}|a_ia_j|\leq m_T$.  Multiplication of Fourier
polynomials is submultiplicative for this coefficient sum; reductions
using $b_i^2=1$ can only decrease it.  Thus the Fourier coefficient
sum of $\Phi_T$ is at most $2^u\sigma^{-p}$.  Lemma~\ref{lem:bias-transfer}
and \eqref{forest:eq:bias} imply
\begin{equation}
 \bigl|\E_{\mathcal Q}\Phi_T-\E_{\mathrm{unif}}\Phi_T\bigr|
 \leq\eta2^u\sigma^{-2(u+1)}
 \leq\frac18\,2^{u-L}\sigma^{2(L-u)}
 \leq\frac18.
 \label{forest:eq:transfer}
\end{equation}
In particular, $\E_{\mathcal Q}\Phi_T\leq1/4$.  The affine formula
$R_p(z)=(p-1+z)/p$ now yields
\[
 \E_{\mathcal Q}C_T
 =\sigma m_T\frac{p-1+\E_{\mathcal Q}\Phi_T}{p}
 \leq\sigma m_T,
\]
as required.

The construction uses only rational sums, products, and integer
powers.  It enumerates $2^u$ assignments and computes powers of
degree $2(u+1)$, so the evaluation time and bit lengths are polynomial
in the input length and $2^u$.
\end{proof}

Different extracted forests may share original labels.  Independence
across forests is neither claimed nor needed: each application of
Lemma~\ref{lem:forest-certificate} bounds an expectation on the same
original sign space, and linearity of expectation bounds any
nonnegative weighted sum of their certificates.  The final algorithm
will minimize such a sum over one common biased space.

% End forests


% Begin odd
\section{Certificates for Odd-Degree XOR}\label{sec:odd}

This section proves the signing theorem for odd degree.  We retain the
original clause labels throughout: two clauses with the same scope still
have different sign variables.  The construction first groups clauses by
common centers, then represents products of residual clauses by a weighted
Kikuchi matrix.  A fractional local condition either bounds the appropriate
memory lift or supplies a forest whose pair products can be certified
separately.  Equalizing the total matrix weight of each original pair makes
this deletion procedure compatible with an exact polynomial identity.

\begin{theorem}[Odd-degree signing certificates]\label{thm:odd-signing}
Fix an odd integer $d\geq3$ and a rational $0<\epsilon<1$.  There is a
constant $C_{d,\epsilon}$ and a deterministic algorithm with the following
property.  Its input consists of $M$ labeled sets $F_i\in\binom{[n]}d$ and
rational coefficients $a_i\in[-1,1]$, where
\[
 M\geq C_{d,\epsilon}n^{d/2}.
\]
The algorithm constructs an explicitly specified inverse-polynomial bias
parameter $\eta>0$ and an exactly evaluable nonnegative rational certificate
$\mathcal C:\{\pm1\}^M\to\mathbb Q_{\geq0}$ such that
\begin{equation}\label{odd:eq:certificate-theorem}
 \sup_{x\in\{\pm1\}^n}
 \left|\sum_{i=1}^M a_i b_i\chi_{F_i}(x)\right|
 \leq\mathcal C(b),\qquad
 \E_{b\sim\mathcal Q}\mathcal C(b)\leq\frac{\epsilon M}{2}
\end{equation}
for every $b$ and every fully $\eta$-biased distribution $\mathcal Q$ on
the original $M$ signs.  Construction and evaluation have polynomial bit
complexity for fixed $d,\epsilon$.
\end{theorem}

We give the proof in several steps.  All constants below depend only on
$d,\epsilon$.  Put
\begin{equation}\label{odd:eq:bundle-constants}
 e=\epsilon/4,\qquad
 B=\lceil64e^{-2}\rceil,\qquad
 R=\lceil128de^{-2}\rceil,\qquad \ell=2d.
\end{equation}
Until the final paragraph, assume $n\geq n_0:=100d\ell$.

\subsection{Bundles and their spread}

For $0\leq q\leq d-1$, define the integer threshold
\begin{equation}\label{odd:eq:bundle-thresholds}
 h_q=\left\lceil BR^q\max\{1,n^{q-d/2}\}\right\rceil.
\end{equation}
Process $q=0,1,\ldots,d-1$ in increasing order.  While a $(d-q)$-set
$C$ is contained in at least $h_q$ currently remaining labels, remove
exactly $h_q$ such labels as a bundle with center $C$.  The residual of
label $i$ in this bundle is $R_i=F_i\setminus C$, of size $q$.  The
algorithm exhausts a stage before proceeding to the next one.  At level
$q$, let $H_q$ denote the number of bundles and let $M_q=H_qh_q$.
Every original label lies in at most one bundle.

\begin{lemma}[Spread and the remainder]\label{odd:lem:spread}
In a level-$q$ bundle, every nonempty residual set $T$ of size $t\leq q$
is contained in at most $h_{q-t}-1$ residuals.  Consequently, for $q>0$,
each residual has at most $qh_{q-1}$ overlapping partners, and
\begin{equation}\label{odd:eq:overlap}
 \frac{qh_{q-1}}{h_q}\leq\frac{2q}{R}
 \leq\eta_{\mathrm{ov}}:=\frac{2d}{R}\leq\frac{e^2}{64}.
\end{equation}
The number $M_{\mathrm{rest}}$ of unbundled labels satisfies
\begin{equation}\label{odd:eq:unbundled}
 M_{\mathrm{rest}}<\frac{nh_{d-1}}d
 \leq U_0n^{d/2},\qquad
 U_0:=\frac{BR^{d-1}+1}{d}.
\end{equation}
\end{lemma}

\begin{proof}
If $h_{q-t}$ labels in the bundle had residuals containing $T$, then all
these labels were still present when the earlier stage $q-t$ terminated.
The set $C\cup T$, of size $d-(q-t)$, would have allowed another bundle
at that stage, a contradiction.  Taking $t=q$ gives the bound
$h_0-1=B-1$ on exact repetitions, so the argument does not require
distinct scopes.

Apply the spread bound to each singleton in a fixed residual and take a
union bound to count overlapping partners.  Since the function
$q\mapsto\max\{1,n^{q-d/2}\}$ is nondecreasing and $BR^q\geq1$,
the ceiling in \eqref{odd:eq:bundle-thresholds} gives
$h_{q-1}/h_q\leq2/R$.  Finally, after the last stage every input vertex
belongs to fewer than $h_{d-1}$ unbundled labels.  Counting incidences
gives $dM_{\mathrm{rest}}<nh_{d-1}$, and
$h_{d-1}\leq BR^{d-1}n^{d/2-1}+1$ proves
\eqref{odd:eq:unbundled}.
\end{proof}

Level $q=0$ consists of bundles of identical full scopes.  Its treatment
will be scalar.  We now fix a positive level, suppress the subscript on
$h=h_q$ and $H=H_q$, and write $M_q=Hh$.

\subsection{Balanced instances and equal total pair weights}

The matrix rows are indexed by $\Omega=\binom{[n]}\ell$, and
$N=|\Omega|$.  Set $a=\lfloor q/2\rfloor$ and $c=\lceil q/2\rceil$;
these letters in this subsection are unrelated to the clause coefficients.
For every unordered within-bundle pair $p=\{i,j\}$ with
$R_i\cap R_j=\varnothing$, include every undirected instance
\begin{equation}\label{odd:eq:balanced-move}
 S\longleftrightarrow S'=S\triangle R_i\triangle R_j,
 \qquad \{|S\cap R_i|,|S\cap R_j|\}=\{a,c\}.
\end{equation}
Both endpoints have size $\ell$.  Each original pair creates exactly
\begin{equation}\label{odd:eq:instance-count}
 T_q=
 \begin{cases}
 \dfrac12\binom q{q/2}^{\!2}\binom{n-2q}{\ell-q},&q\text{ even},\\[4pt]
 \binom q{\lfloor q/2\rfloor}^{\!2}\binom{n-2q}{\ell-q},&q\text{ odd}
 \end{cases}
\end{equation}
instances.  Indeed, an even $q$ permits one intersection pattern and an
odd $q$ permits two, after which each undirected instance has been counted
twice.  The instances remain labeled by their original pair even when
different pairs have identical endpoints.
For a nonempty positive level, the pair collection is nonempty:
\eqref{odd:eq:overlap} and $1/h+\eta_{\mathrm{ov}}<1$ ensure that every
bundle has a disjoint pair.

We use the reciprocal-excess weights of
Schmidhuber--Hastings~\cite[Section~3.2]{SH26Moore}.  An instance $f$ with
endpoints $S,S'$ and labels $i,j$ has four ports
\[
 \mathcal P(f)=\{(S,i),(S,j),(S',i),(S',j)\}.
\]
Let $c(z)$ be the number of initial instances using port $z$, and put
\begin{equation}\label{odd:eq:port-weight}
 \delta_f=1+\sum_{z\in\mathcal P(f)}(c(z)-1),\qquad
 w_f=\delta_f^{-1}.
\end{equation}
Since $\delta_f\geq c(z)$ for every port of $f$,
\begin{equation}\label{odd:eq:port-capacity}
 \sum_{f:\,z\in\mathcal P(f)}w_f\leq1.
\end{equation}

\begin{lemma}[Per-pair contention and equalization]\label{odd:lem:equalize}
There is an explicit constant
\[
 \Xi:=16BR^{d-1}(8\ell)^{d-1}
\]
such that, for every original disjoint pair $p$,
\begin{equation}\label{odd:eq:pair-contention}
 \frac1{T_q}\sum_{f\text{ of }p}(\delta_f-1)\leq\Xi.
\end{equation}
One can construct positive rational weights $\widetilde w_f\leq w_f$
and a common rational $\beta>0$ such that
\begin{equation}\label{odd:eq:equal-pair-mass}
 \sum_{f\text{ of }p}\widetilde w_f=\beta
 \quad\text{for every }p,\qquad
 \beta\geq\frac{T_q}{1+\Xi}.
\end{equation}
Their port loads remain at most one.  With $K_0:=2(1+\Xi)$, they satisfy
\begin{equation}\label{odd:eq:beta-estimate}
 \frac N\beta\leq K_0n^q,\qquad
 \frac{N}{\beta h}\leq K_0n^{d/2}.
\end{equation}
\end{lemma}

\begin{proof}
Fix $p=\{i,j\}$.  Choose one of its instances uniformly, then choose one
of that instance's two endpoints uniformly and call it $S$.  This is
equivalent to choosing a balanced starting row uniformly.  We first bound
the expected excess at the port $(S,i)$.
Any competing instance at this port has partner $j'$ from the same
bundle as $i$, with $R_{j'}\cap R_i=\varnothing$.  A specified partner
determines at most one instance at row $S$.  Put
$t=|R_{j'}\cap R_j|$.  For a specified trace $T\subseteq R_j$ of size
$t>0$, the number of possible partners is at most $h_{q-t}$, by
Lemma~\ref{odd:lem:spread}; for $t=0$ it is at most $h_q$.

Condition on $S\cap R_i$ and $S\cap R_j$.  The other $\ell-q$ row
vertices form a uniformly random subset of the $n-2q$ vertices outside
$R_i\cup R_j$.  If $t<a$, an eligible $j'$ requires at least $a-t$
vertices of $R_{j'}\setminus R_j$ in this random subset.  Its probability
is therefore at most
\begin{equation}\label{odd:eq:contention-probability}
 \binom{q-t}{a-t}
 \left(\frac{\ell-q}{n-3q}\right)^{a-t}.
\end{equation}
This follows by a union bound over the chosen $a-t$ vertices and the
falling-factorial probability of including them.  Since $n-3q\geq n/2$,
the product of \eqref{odd:eq:contention-probability} and the partner bound
is at most
\[
 2BR^q\,2^q(2\ell)^q
 \bigl(n^{-(a-t)}+n^{c-d/2}\bigr)
 \leq4BR^q\,2^q(2\ell)^q.
\]
Here the decisive inequality is
$c=\lceil q/2\rceil\leq(d-1)/2<d/2$.  If $t\geq a$, use probability
one instead: then $q-t\leq c<d/2$, so $h_{q-t}=BR^{q-t}\leq BR^q$.
This includes $t=q$, and hence includes repeated residuals.  Sum over at
most $2^q$ traces and repeat the estimate after exchanging $i,j$.
Because $S$ was a random endpoint, twice these two expected port excesses
equal the average summed excess at all four ports of an undirected
instance.  The resulting bound is
$16BR^q(8\ell)^q\leq\Xi$.  Averaging over the one or two possible
intersection patterns does not change the estimate.  This proves
\eqref{odd:eq:pair-contention}.

Write $W_p=\sum_{f\text{ of }p}w_f$.  Cauchy--Schwarz gives
\[
 W_p\geq\frac{T_q^2}{\sum_{f\text{ of }p}\delta_f}
 \geq\frac{T_q}{1+\Xi}.
\]
Take $\beta=\min_pW_p$ and set
$\widetilde w_f=(\beta/W_p)w_f$ for each instance of $p$.
The weights decrease and have exactly the required pair sums.

To verify the dimension bound, use
\[
 \frac{\binom{n-2q}{\ell-q}}{\binom n\ell}
 =\frac{(\ell)_q(n-\ell)_q}{(n)_{2q}}
 \geq\left(\frac{\ell}{4n}\right)^q.
\]
The lower bound holds because $\ell\geq2q$ and
$n-\ell-q\geq n/2$.  The coefficient multiplying this ratio in
\eqref{odd:eq:instance-count} is at least $1/2$.  Thus
$N/\beta\leq2(1+\Xi)(4n/\ell)^q\leq K_0n^q$, since $\ell\geq6$.
Finally $h\geq\max\{1,n^{q-d/2}\}$ gives the second inequality in
\eqref{odd:eq:beta-estimate}.
\end{proof}

We freeze $\beta$ and all the equalized weights.  Every subsequent deletion
removes an entire original pair and all of its instances.  It changes no
other pair weight.  Thus both equalization and port capacity survive every
iteration.

\subsection{The comparison after pair deletion}

Let $E$ be any remaining collection of disjoint within-bundle pairs.  Let
$\mathcal T$ be any edge-disjoint family of forests whose edges, together
with $E$, partition the initial disjoint pairs.  These are forests on
original clause labels; different forests may share labels.  Define
\[
 Q_T(b,x)=\sum_{\{i,j\}\in T}a_i a_jb_i b_j
                         \chi_{R_i\triangle R_j}(x),
 \qquad
 P_q(b,x)=\sum_{\mathcal B}\chi_{C_{\mathcal B}}(x)
                       \sum_{i\in\mathcal B}a_i b_i\chi_{R_i}(x).
\]
Let $A_b$ be the symmetric matrix that assigns weight
$\widetilde w_f a_i a_jb_i b_j$ to each surviving instance of $\{i,j\}$.
Let $A_+$ use weight $\widetilde w_f$, ignoring the coefficients as well
as the signs.  For $E\ne\varnothing$, put
\begin{equation}\label{odd:eq:normalization}
 d(S)=\sum_{S'}A_+[S,S'],\qquad
 \overline d=\frac{2\beta|E|}{N},\qquad
 \Gamma=\operatorname{diag}(d(S))+\overline d I,\qquad
 K_b=\Gamma^{-1/2}A_b\Gamma^{-1/2}.
\end{equation}

\begin{lemma}[Comparison with the original level polynomial]
\label{odd:lem:comparison}
With the preceding notation,
\begin{equation}\label{odd:eq:comparison}
 \frac{\sup_x|P_q(b,x)|^2}{M_q^2}
 \leq\frac1h+\eta_{\mathrm{ov}}
   +\frac{2}{M_qh}\sum_{T\in\mathcal T}\sup_x|Q_T(b,x)|
   +2\|K_b\|.
\end{equation}
The last term may instead be replaced by $2|E|/(M_qh)$, including the
case $E=\varnothing$ without forming a normalized matrix.
\end{lemma}

\begin{proof}
For the row vector $z(x)_S=\chi_S(x)$, equalization gives the exact identity
\begin{equation}\label{odd:eq:polynomial-identity}
 z(x)^{\mathsf T}A_bz(x)
 =2\beta\sum_{\{i,j\}\in E}a_i a_jb_i b_j
                                      \chi_{R_i\triangle R_j}(x).
\end{equation}
Let $O$ denote the initial overlapping pairs.  Apply Cauchy--Schwarz once,
using the original $H$ bundles, to obtain
\begin{align*}
 |P_q(b,x)|^2
 &\leq H\sum_{\mathcal B}
              \left(\sum_{i\in\mathcal B}a_i b_i\chi_{R_i}(x)\right)^2\\
 &=H\sum_i a_i^2
   +2H\sum_{\{i,j\}\in O}a_i a_jb_i b_j\chi_{R_i\triangle R_j}(x)
   +2H\sum_{T\in\mathcal T}Q_T(b,x)
   +\frac H\beta z(x)^{\mathsf T}A_bz(x).
\end{align*}
The first term is at most $HM_q=M_q^2/h$.  By
\eqref{odd:eq:overlap}, $|O|\leq\eta_{\mathrm{ov}}M_qh/2$.
Moreover,
\[
 |z(x)^{\mathsf T}A_bz(x)|
 \leq\|K_b\|\,z(x)^{\mathsf T}\Gamma z(x)
 =4\beta|E|\|K_b\|.
\]
Use $|E|\leq H\binom h2\leq M_qh/2$ and divide by $M_q^2$.
Paying each surviving pair by one gives the alternate last term directly.
In particular, the number of deleted forests does not affect the
Cauchy--Schwarz factor or create additional diagonal terms.
\end{proof}

We will stop a level if its current average degree is small.  For any
threshold $d_*>0$, \eqref{odd:eq:beta-estimate} gives
\begin{equation}\label{odd:eq:low-degree}
 \overline d<d_*\quad\Longrightarrow\quad
 \frac{2|E|}{M_qh}
 =\frac{N\overline d}{\beta M_qh}
 \leq K_0d_*\frac{n^{d/2}}{M_q}.
\end{equation}
Thus it suffices to retain only levels with a sufficiently large original
count $M_q$; there are at most $d-1$ other positive levels to discard.

\subsection{A fractional window condition and the memory lift}

For a set $U\subseteq[n]$, take the graph on labels whose edges are the
surviving pairs $\{i,j\}$ with $R_i,R_j\subseteq U$.  At capacity $D$,
call $U$ \emph{good} if the following linear system is feasible:
\begin{equation}\label{odd:eq:window-lp}
 \lambda_{i,v}\geq0\quad(v\in R_i),\qquad
 t_i:=\sum_{v\in R_i}\lambda_{i,v},\qquad
 \sum_i\lambda_{i,v}\leq D\quad(v\in U),\qquad
 t_i+t_j\geq1\quad(\{i,j\}\in E[U]).
\end{equation}
One may impose $t_i\leq1$ by truncating each larger vector proportionally;
this preserves all pair inequalities and decreases all vertex loads.
Testing goodness is a rational linear program of polynomial size.

The original-sign parity lift has states $(S,\mu)$, where
$S\in\Omega$ and $\mu\subseteq[M]$ records the active original labels.
An instance of pair $\{i,j\}$ changes the state by
\[
 (S,\mu)\longleftrightarrow
 (S\triangle R_i\triangle R_j,\,\mu\triangle\{i,j\}).
\]
Its unnormalized weight is $\widetilde w_f$, and its normalized weight
divides this by $\sqrt{\Gamma(S)\Gamma(S')}$.
Associate to a state its residual window
\[
 W(S,\mu)=S\cup\bigcup_{i\in\mu}R_i.
\]
An empty-memory closed walk of length $2s$ has at most
$\min\{2j,2(2s-j)\}\leq2s$ active labels after $j$ steps.  All its
windows therefore have size at most
\begin{equation}\label{odd:eq:window-size}
 L=\ell+2ds.
\end{equation}

\begin{lemma}[Norm on good windows]\label{odd:lem:good-norm}
Restrict the normalized parity lift to states with $|\mu|\leq2s$ and
good windows.  Its adjacency operator $\mathcal L_{\mathrm{good}}$ satisfies
\begin{equation}\label{odd:eq:good-norm}
 \|\mathcal L_{\mathrm{good}}\|
 \leq2\sqrt{\frac{(D+1)\ell}{\overline d}}.
\end{equation}
\end{lemma}

\begin{proof}
Fix one feasible solution to \eqref{odd:eq:window-lp} for each good
window.  These solutions are used only to prove the bound; the algorithm
does not enumerate all windows.  We split and orient each lifted edge
so that the total incoming unnormalized weight is at most $(D+1)\ell$.

First consider an edge whose two endpoint windows equal $U$, with pair
$\{i,j\}$.  Both residuals lie in $U$.  Split its weight among vertices
of these two residuals in fractions
\[
 \frac{\lambda_{i,v}}{t_i+t_j}\quad(v\in R_i),\qquad
 \frac{\lambda_{j,v}}{t_i+t_j}\quad(v\in R_j).
\]
The fractions sum to one.  Since the residuals are disjoint and their
union is $S\triangle S'$, each chosen vertex belongs to exactly one
endpoint row.  Orient the corresponding piece toward that row.
At a fixed state with row $S$, the incoming weight witnessed by $v\in S$
is at most $\sum_i\lambda_{i,v}\leq D$: for a fixed $i$ the relevant
instances use the port $(S,i)$, whose total weight is at most one, and
$t_i+t_j\geq1$.  Thus equal-window edges contribute at most $D\ell$.

For an unequal-window edge, choose an endpoint $x=(S,\mu)$ whose window
contains a vertex absent from the other endpoint's window, breaking ties
deterministically.  Orient toward $x$ and choose such a vertex $v$ as
witness.  Write the other endpoint as $y=(S',\mu')$.  The witness belongs
to the residual of exactly one toggled label, say $i$, because the two
residuals are disjoint.  It belongs to $S$: otherwise the row toggle
would put it in $S'\subseteq W(y)$.  Label $i$ was entered, since a
removed label would put its residual in $W(y)$.  Furthermore, $i$ is the
unique label in $\mu$ whose residual contains $v$.  Any other such label
would either remain in $\mu'$ or be the other toggled label, and both
possibilities contradict the choice of $v$.  Thus $(x,v)$ determines
$i$, and all incoming edges with witness $v$ use the port $(S,i)$.
They have total weight at most one.  There are at most $\ell$ witnesses.
This is the unequal-window witness argument of
\cite[Section~3.3]{SH26Moore}, applied to labeled residuals.

The lifted degree at a state with row $S$ is at most $d(S)$.
Lemma~\ref{lem:orientation}, with normalization
$\Gamma(S)=d(S)+\overline d$, now gives
\eqref{odd:eq:good-norm}.  Fractionally splitting an edge is harmless for
that lemma: each piece has nonnegative weight and the pieces sum to the
original edge.
\end{proof}

If $U$ is not good, Lemma~\ref{lem:forest-alternative} constructs a
forest $T$ of surviving pairs and its residual union $U_T\subseteq U$
such that
\begin{equation}\label{odd:eq:dense-forest}
 |T|>D|U_T|.
\end{equation}
All residuals are nonempty because $q>0$.  This alternative applies to
the current pair graph, with no requirement that it remain a union of
bundle cliques.  It is the reason whole-pair deletion can be repeated.

\subsection{Deterministic decomposition by a trace test}

Set
\begin{equation}\label{odd:eq:trace-constants}
 \begin{gathered}
 \sigma=\frac{e^2}{64},\qquad \zeta=\frac{e^2}{128},\qquad
 D=\lceil32\sigma^{-2}\rceil,\\
 s=\lceil\log_2(128N)\rceil,\qquad
 \beta_0=\zeta^{2s},\qquad
 d_* =\frac{16(D+1)\ell}{\zeta^2}.
 \end{gathered}
\end{equation}
Here $\beta_0$ is a trace threshold, distinct from the common pair mass
$\beta$.  Set the positive rational bias parameter to
\begin{equation}\label{odd:eq:bias}
 \eta=\min\left\{
 \frac{\beta_0}{512N},\
 \frac{\sigma^{2(L+1)}}{8\cdot2^L},\
 \frac{e^2}{8}\right\},\qquad L=\ell+2ds.
\end{equation}
For fixed $d,e$, its reciprocal is polynomial in $n$.
Finally set
\begin{equation}\label{odd:eq:retained-threshold}
 A=\left\lceil\frac{64K_0d_*}{e^2}\right\rceil
\end{equation}
and retain a positive level only if $M_q\geq An^{d/2}$.  By
\eqref{odd:eq:low-degree}, a low-degree stop on a retained level costs
at most $e^2/64$ in the squared comparison.

Fix one explicit fully $\eta$-biased distribution $\mathcal Q_0$ on the
original $M$ signs, using Lemma~\ref{lem:small-bias}.  On each retained
level start with all its disjoint pairs and repeat:
\begin{enumerate}
\item If $E=\varnothing$ or $\overline d<d_*$, stop and use the direct
      pair-count bound.
\item Form the normalized matrix $K_b^{(0)}$ using instance weights
      $\widetilde w_fb_ib_j$, without the $a_i a_j$ coefficients.  If
      \begin{equation}\label{odd:eq:trace-stop}
       \E_{b\sim\mathcal Q_0}\operatorname{tr}
                   (K_b^{(0)})^{2s}\leq\beta_0/32,
      \end{equation}
      stop and retain this trace test.
\item Otherwise find a failed window, extract a forest by
      \eqref{odd:eq:dense-forest}, save that forest, delete all its
      pairs and their instances, and repeat.
\end{enumerate}

We next justify constructively the third step.  The matrix
$P_b=\Gamma^{-1}A_b^{(0)}$ is similar to $K_b^{(0)}$, so its traces are
identical and rational.  Every Fourier coefficient in an entry of
$P_b^j$ is nonnegative.  Its coefficient mass is at most one, because
the unsigned matrix $P_+=\Gamma^{-1}A_+$ has row sums at most one.
Likewise the nonnegative Fourier coefficient mass of
$\operatorname{tr}(K_b^{(0)})^{2s}$ is
$\operatorname{tr}(P_+^{2s})\leq N$.  Small-bias transfer therefore changes
its expectation by at most $\eta N$.

Under uniform original signs, the surviving trace terms are precisely
the empty-memory closed walks of length $2s$.  Failure of
\eqref{odd:eq:trace-stop} implies that their total mass is greater than
\[
 \beta_0/32-\eta N\geq15\beta_0/512.
\]
Since $\overline d\geq d_*$, Lemma~\ref{odd:lem:good-norm} bounds the
mass of those walks staying in good states by
\[
 N\|\mathcal L_{\mathrm{good}}\|^{2s}
 \leq N(\zeta^2/4)^s
 =\beta_0 N4^{-s}\leq\beta_0/512.
\]
Here there are $N$ empty-memory starting rows, and their length-$2s$
diagonal entries are individually bounded by the operator norm to that
power.  If a starting window itself is bad, walks starting there simply
belong to the bad part.  Thus the total mass of bad empty-memory closed
walks is greater than $\delta:=\beta_0/64$.

Apply Theorem~\ref{thm:extraction} to the substochastic matrix $P_b$,
with the test \eqref{odd:eq:window-lp} on recovered memories.  It recovers
a bad state using threshold
\begin{equation}\label{odd:eq:extraction-threshold}
 \tau=\frac{\delta}{2N(2s+1)}.
\end{equation}
Only memories with at most $2s$ labels are tested, as permitted by the
empty-memory closed-walk bound.  Every returned window therefore obeys
$|U|\leq L$.  Its inverse threshold is polynomial in $n$, so this step
is deterministic and polynomial time.  Lemma~\ref{lem:forest-alternative}
then supplies the desired forest.  Each iteration removes at least one
original pair, and there are at most $H\binom h2\leq M_qh/2$ initial
pairs.  The process terminates after polynomially many iterations.

The stopping test also controls arbitrary coefficients and any final
small-biased distribution.  Write $K_b^{(a)}$ for the normalized matrix
with the actual $a_i a_j$ coefficients.  In a uniform-sign trace
expansion, every surviving term contains each original label an even
number of times.  Its coefficient factor has the form
$\prod_i a_i^{2r_i}\in[0,1]$.  Consequently
\[
 \E_{\mathrm{unif}}\operatorname{tr}(K_b^{(a)})^{2s}
 \leq\E_{\mathrm{unif}}\operatorname{tr}(K_b^{(0)})^{2s}.
\]
The absolute Fourier coefficient mass of the actual trace is at most
$N$, by domination by the unsigned path expansion.  For every fully
$\eta$-biased $\mathcal Q$, a trace stop thus implies
\begin{equation}\label{odd:eq:actual-trace}
 \E_{b\sim\mathcal Q}\operatorname{tr}(K_b^{(a)})^{2s}
 \leq\frac{\beta_0}{32}+2\eta N<\frac{\beta_0}{8}.
\end{equation}
This comparison treats negative and zero coefficients without changing
the pair weights, the density test, or the decomposition.

\subsection{Rational certificates and the final expectation}

For each extracted forest $T$, let $m_T=|T|$ and $u=|U_T|\leq L$.
The products $b_ib_j$ on its edges are independent uniform signs under
uniform original signs.  Lemma~\ref{lem:forest-certificate} applies to
$Q_T$ with the parameters in \eqref{odd:eq:trace-constants} and
\eqref{odd:eq:bias}.  Explicitly, its rational potential and certificate
are
\begin{equation}\label{odd:eq:forest-certificate}
 p_T=2(u+1),\qquad
 \Phi_T(b)=\sum_{x\in\{\pm1\}^{U_T}}
          \left(\frac{Q_T(b,x)}{\sigma m_T}\right)^{p_T},\qquad
 C_T(b)=\sigma m_T R_{p_T}(\Phi_T(b)),
\end{equation}
where $R_p$ is the rational envelope in \eqref{eq:young-envelope}.
They satisfy, pointwise and for every fully $\eta$-biased $\mathcal Q$,
\begin{equation}\label{odd:eq:forest-expectation}
 C_T(b)\geq\sup_x|Q_T(b,x)|,\qquad
 \E_{\mathcal Q}C_T(b)\leq\sigma m_T.
\end{equation}
The degree condition $m_T>Du$ is exactly
\eqref{odd:eq:dense-forest}; different forests may reuse original
labels, and \eqref{odd:eq:forest-expectation} does not assume independence
between different forests.

At a trace stop define
\begin{equation}\label{odd:eq:matrix-certificate}
 \Psi_q(b)=\frac{\operatorname{tr}(K_b^{(a)})^{2s}}{\zeta^{2s}},
 \qquad C_{K,q}(b)=\zeta R_{2s}(\Psi_q(b)).
\end{equation}
The matrix is real symmetric, so its even trace is nonnegative and
dominates its operator norm to the power $2s$.  The envelope and
\eqref{odd:eq:actual-trace} give
\begin{equation}\label{odd:eq:matrix-expectation}
 C_{K,q}(b)\geq\|K_b^{(a)}\|,\qquad
 \E_{\mathcal Q}C_{K,q}(b)\leq\zeta.
\end{equation}
The trace can be computed as
$\operatorname{tr}(\Gamma^{-1}A_b)^{2s}$; no irrational square roots
are required in the algorithm.

In \eqref{odd:eq:comparison}, replace each forest maximum by $C_T$.
Use $2C_{K,q}$ at a trace stop, and $2|E|/(M_qh)$ at a low-degree stop.
Call the resulting rational bound $B_q(b)$.  Because the deleted forests
are edge-disjoint,
\[
 \sum_{T\in\mathcal T}m_T\leq M_qh/2.
\]
Together with \eqref{odd:eq:overlap},
\eqref{odd:eq:retained-threshold}, and
\eqref{odd:eq:forest-expectation}--\eqref{odd:eq:matrix-expectation},
this gives
\begin{equation}\label{odd:eq:squared-budget}
 \E_{\mathcal Q} B_q
 \leq\frac{e^2}{64}+\frac{e^2}{64}+\sigma+
                      \frac{e^2}{64}
 =\frac{e^2}{16}.
\end{equation}
The four terms pay respectively for the original diagonal, overlapping
pairs, all deleted forests, and the final remainder.  The last term is
valid for both kinds of stop, since $2\zeta=e^2/64$.

Using the square-root envelope \eqref{eq:sqrt-envelope}, define
\begin{equation}\label{odd:eq:level-certificate}
 \mathcal C_q(b)=M_q\left(\frac e2+\frac{B_q(b)}{2e}\right).
\end{equation}
It is pointwise at least $\sup_x|P_q(b,x)|$, and
\begin{equation}\label{odd:eq:level-expectation}
 \E_{\mathcal Q}\mathcal C_q
 \leq\frac{17}{32}eM_q<eM_q.
\end{equation}
This uses an aggregate expectation for all forests rather than a union
bound over the number of extractions.

For a level-zero bundle $\mathcal B$, put
$V_{\mathcal B}(b)=\sum_{i\in\mathcal B}a_i b_i$.  Its maximum absolute
contribution is exactly $|V_{\mathcal B}(b)|$.  Use
\[
 \Phi_{\mathcal B}(b)=\left(\frac{V_{\mathcal B}(b)}{eB}\right)^2,
 \qquad
 C_{\mathcal B}(b)=eB R_2(\Phi_{\mathcal B}(b)).
\]
Its pointwise soundness follows from the envelope, and small bias gives
\[
 \E_{\mathcal Q}\Phi_{\mathcal B}
 \leq\frac{B+\eta B^2}{e^2B^2}
 \leq\frac1{64}+\frac18<\frac14.
\]
Hence $\E_{\mathcal Q}C_{\mathcal B}\leq eB$.

Let $M_{\mathrm{err}}$ be the total number of unbundled labels and labels
in discarded positive levels.  Equations~\eqref{odd:eq:unbundled} and
\eqref{odd:eq:retained-threshold} imply
\[
 M_{\mathrm{err}}\leq\bigl(U_0+(d-1)A\bigr)n^{d/2}.
\]
Choose
\begin{equation}\label{odd:eq:final-density-constant}
 C_{d,\epsilon}\geq
 \max\left\{
 \frac{U_0+(d-1)A}{e},\
 16e^{-2}(n_0+1)\right\}.
\end{equation}
Then $M_{\mathrm{err}}\leq eM$.  Sum its trivial count, all level-zero
certificates, and all retained level certificates to define $\mathcal C$.
The triangle inequality gives pointwise soundness, and the preceding
expectations give
\[
 \E_{\mathcal Q}\mathcal C
 \leq M_{\mathrm{err}}+e\sum_{q=0}^{d-1}M_q
 \leq2eM=\epsilon M/2.
\]
All terms are evaluated on the same original sign vector.  In particular,
the argument permits adding certificates constructed independently for
different Fourier layers.

\paragraph{Small input lengths.}
If $n<n_0$, use the original polynomial
$P(b,x)=\sum_i a_i b_i\chi_{F_i}(x)$ directly.  Set
\[
 p=2(n+1),\qquad
 \Phi(b)=\sum_{x\in\{\pm1\}^n}
                   \left(\frac{P(b,x)}{eM}\right)^p,
 \qquad
 \eta=\frac{e^{2(n+1)}}{8\cdot2^n},
 \qquad \mathcal C(b)=eM R_p(\Phi(b)).
\]
The density choice \eqref{odd:eq:final-density-constant} ensures
$M\geq16e^{-2}(n+1)$.  By Lemma~\ref{lem:moment},
\[
 \E_{\mathrm{unif}}\Phi
 \leq2^n\left(\frac{2(n+1)}{e^2M}\right)^{n+1}
 \leq2^n8^{-(n+1)}\leq\frac18.
\]
The absolute Fourier coefficient mass of $\Phi$ is at most
$2^ne^{-p}$, so the displayed bias bounds the transfer error by $1/8$.
Thus $\E_{\mathcal Q}\Phi\leq1/4$ and
$\E_{\mathcal Q}\mathcal C\leq eM<\epsilon M/2$.
Pointwise soundness again follows from the even moment and its rational
envelope.  Enumerating $2^n$ inputs has constant cost depending on
$d,\epsilon$ in this branch; it does not enumerate the $2^M$ signings.

\paragraph{Bit complexity.}
At fixed $d,\epsilon$, there are polynomially many bundles, pairs, rows,
and pair instances.  All initial port loads are integers of polynomial
size.  The rational sums $W_p$ and the equalized weights therefore have
polynomial bit length.  They are frozen during the deletion loop, so
successive iterations do not recursively compound their denominators.
Each trace uses a rational matrix power of order $O(\log n)$; each local
certificate enumerates $2^{O(\log n)}$ assignments and uses moments of
order $O(\log n)$.  The explicit small-biased spaces, recovered memory
lists, rational linear programs, and forest extractions all have
polynomial size.  Integer square-root comparisons compute the ceilings
in \eqref{odd:eq:bundle-thresholds} exactly.  These observations establish
the construction and evaluation claims in Theorem~\ref{thm:odd-signing}.

% End odd


% Begin avoidance
\section{From Signing Certificates to Range Avoidance}
\label{sec:avoidance}

We first assemble the degree-$d$ certificate, including the conversion
from coefficient mass to label count. We then apply the Fourier
reduction of Guruswami, Lyu, and Yuan~\cite[Section~6.1]{GLY26}.
The proof selects one sign vector for all Fourier layers.

\subsection{The Signing Theorem and the Elementary Degrees}

\begin{proof}[Proof of \Cref{thm:signing}]
For odd $d\ge3$, apply \Cref{thm:odd-signing}.
Suppose $d$ is even, and write $W=\sum_i|a_i|\le M$.
If $W\le\epsilon M/2$, the constant certificate
$\mathcal C(b)=W$ has both required properties; one may choose any
rational $\eta\in(0,1)$. Otherwise $W>\epsilon M/2$.
Let $C^{\mathrm{wt}}_{d,\epsilon}$ denote a sufficient density
constant in \Cref{thm:even-weight}. Choose the label-count constant
so that
\[
 C_{d,\epsilon}\ge(2/\epsilon)C^{\mathrm{wt}}_{d,\epsilon}.
\]
The density hypothesis then gives
$W>C^{\mathrm{wt}}_{d,\epsilon}n^{d/2}$. The even-degree theorem
constructs a certificate with
\[
 \E\mathcal C(b)\le\frac{41}{128}\epsilon W
 \le\frac{41}{128}\epsilon M<\epsilon M/2.
\]
Its pointwise guarantee and bit complexity are exactly those needed.
The bias parameter in either case can be decreased, if necessary, to
a rational inverse-polynomial value shared with other systems.
\end{proof}

The degrees zero and one need a different density scale from
$n^{d/2}$. Their certificates are elementary.

\begin{lemma}[Degree zero and degree one]
\label{lem:elementary-degrees}
Let $n\ge1$, let $a_i\in\Q\cap[-1,1]$, and suppose either
all $F_i$ are singletons or all $F_i$ are empty. Define
\begin{equation}
 \mathcal C_1(b)=\sum_{v=1}^n
       \abs{\sum_{i:F_i=\{v\}}a_ib_i},
 \qquad
 \mathcal C_0(b)=\abs{\sum_i a_ib_i},
 \label{eq:elementary-certificates}
\end{equation}
in the respective cases. These are exact certificates for the
maximum absolute value. Under every fully $\eta$-biased
distribution,
\begin{equation}
 \E\mathcal C_1\le\sqrt{nM}+\sqrt\eta M,
 \qquad
 \E\mathcal C_0\le\sqrt M+\sqrt\eta M.
 \label{eq:elementary-expectation}
\end{equation}
In particular, if $M\ge16\epsilon^{-2}n$ and
$\eta\le\epsilon^2/16$, either expectation is at most
$\epsilon M/2$.
\end{lemma}
\begin{proof}
For degree one, each variable can be chosen to align with its grouped
coefficient, proving exactness. The degree-zero assertion is immediate.
For a group $I$ of original labels, small bias bounds its second
moment by
\[
 \E\left(\sum_{i\in I}a_ib_i\right)^2
 \le\sum_{i\in I}a_i^2+
       \eta\left(\sum_{i\in I}|a_i|\right)^2.
\]
Apply Cauchy--Schwarz, the inequality $\sqrt{A+B}\le\sqrt A+\sqrt B$,
and then Cauchy--Schwarz across at most $n$ disjoint groups.
Since $\sum_i a_i^2\le M$ and $\sum_i|a_i|\le M$, this gives
Eqn.~\eqref{eq:elementary-expectation}. The stated parameters make
each of its two summands at most $\epsilon M/4$.
\end{proof}

\subsection{Removing Parities and the Top Fourier Term}

\begin{lemma}[The Fourier gap]
\label{lem:fourier-gap}
Fix $k\ge3$ and put $\delta=2^{1-k}$. Let
$g:\{\pm1\}^k\to\{\pm1\}$ be a Boolean function that is
not a signed parity. There is a rational multilinear polynomial
$q$ of degree at most $k-1$ such that
\begin{equation}
 |g(x)-q(x)|\le1-\delta\quad\text{for every }x.
 \label{eq:fourier-gap}
\end{equation}
Each coefficient of $q$ has magnitude at most one, and its Fourier
expansion is computable in $O_k(1)$ time.
\end{lemma}
\begin{proof}
Write $g=\sum_{A\subseteq[k]}\widehat g(A)\chi_A$ and put
$a=\widehat g([k])$. The sign function $g\chi_{[k]}$ is
nonconstant: otherwise $g$ would be a signed parity. Its mean
therefore lies in $[-1+2^{1-k},1-2^{1-k}]$. Taking
$q=g-a\chi_{[k]}$ proves Eqn.~\eqref{eq:fourier-gap}.
The remaining claims follow from
$\widehat g(A)=2^{-k}\sum_xg(x)\chi_A(x)$.
\end{proof}

If a coordinate depends on fewer than $k$ inputs and $n\ge k$,
pad its neighborhood to $k$ distinct variables and extend its truth
table by ignoring the added variables. This preserves the function.
It permits every coordinate to use the same set of local-position
labels $[k]$, even when their neighborhoods differ.

\begin{proof}[Proof of \Cref{thm:avoid}]
Increase $C_k$ so that the promise implies $m>n$. If $n<k$,
enumerate all $2^n$ inputs and their images. A missing string can
be found among the first $2^n+1$ strings of length $m$, in time
$O_k(m)$ with a suitable dictionary. Thus assume $n\ge k$.

Identify every coordinate that is a signed parity; constant
coordinates are included as parities of the empty set.
Suppose there are $p>n$ such coordinates, written as
$f_i(x)=s_i\chi_{T_i}(x)$, $s_i\in\{\pm1\}$. Their incidence
vectors have a nontrivial binary dependency $\lambda\in\F_2^p$.
Consequently every output of this submap satisfies
\[
 \prod_{i:\lambda_i=1}f_i(x)
 =\prod_{i:\lambda_i=1}s_i.
\]
Choose target signs violating this equality and extend them arbitrarily
to the other outputs. Binary linear algebra finds the dependency and
the violating target in polynomial time.

Otherwise discard these $p\le n$ coordinates and retain $M=m-p$
nonparity outputs. Pad and order each remaining input neighborhood as
$v_{i,1},\ldots,v_{i,k}$. Apply \Cref{lem:fourier-gap} to obtain
\begin{equation}
 f_i(x)=\widehat f_i([k])\chi_{\{v_{i,1},\ldots,v_{i,k}\}}(x)
            +q_i(x),\qquad
 |f_i(x)-q_i(x)|\le1-\delta.
 \label{eq:coordinate-split}
\end{equation}
For every proper subset $A\subsetneq[k]$, define one layer by
\[
 F_{i,A}=\{v_{i,j}:j\in A\},\qquad
 P_A(b,x)=\sum_{i=1}^M\widehat f_i(A)b_i\chi_{F_{i,A}}(x).
\]
There are $J=2^k-1$ layers and
$\sum_i b_iq_i(x)=\sum_{A\subsetneq[k]}P_A(b,x)$.
Zero Fourier coefficients remain as zero-weight original labels.
In particular, every layer has exactly $M$ labels on the same
original sign space.

Set $\epsilon_0=\delta/(4J)$. For $|A|\ge2$, apply
\Cref{thm:signing} with parameter $\epsilon_0$. For $|A|\le1$,
use \Cref{lem:elementary-degrees}. Taking $C_k$ sufficiently large
ensures all density conditions: for $2\le d\le k-1$,
$n^{d/2}\le n^{(k-1)/2}$, and
$n\le n^{(k-1)/2}$. The at most $n$ discarded coordinates are
absorbed by increasing $C_k$ once more.

Let $\eta_*$ be the minimum of the finitely many required bias
parameters and $\epsilon_0^2/16$. Construct one explicit fully
$\eta_*$-biased space $\mathcal Q_*$ on the $M$ original output
signs. Each layer supplies a pointwise certificate $\mathcal C_A$,
and their sum satisfies
\begin{equation}
 \E_{b\sim\mathcal Q_*}\sum_{A\subsetneq[k]}\mathcal C_A(b)
 \le J\epsilon_0M/2=\delta M/8.
 \label{eq:joint-layer-budget}
\end{equation}
Enumerate this space and select a sign vector whose summed
certificate is at most $\delta M/8$. In particular,
\begin{equation}
 \sup_x\abs{\sum_i b_iq_i(x)}<\delta M.
 \label{eq:final-separation}
\end{equation}
If an input mapped to this target on all retained coordinates, then
\[
 \sum_i b_iq_i(x)
 \ge\sum_i b_if_i(x)-\sum_i|f_i(x)-q_i(x)|
 \ge M-(1-\delta)M=\delta M,
\]
contradicting Eqn.~\eqref{eq:final-separation}.
Extend $b$ arbitrarily to the discarded positions and convert back
to bits. The resulting string is outside the full range.

All truth-table operations have constant cost per coordinate for
fixed $k$. The number of layers is constant. Each certificate is
constructed and evaluated in polynomial bit complexity, and
$\eta_*^{-1}$ is polynomial. Enumeration of $\mathcal Q_*$ is
therefore polynomial in the explicit input length. This proves
the deterministic running-time bound.
\end{proof}

For clarity, the algorithm implemented by the proof is the following.
\begin{enumerate}
\item Handle bounded input size by exhaustive input evaluation.
Identify signed-parity coordinates and, if there are more than $n$,
return a target violating one binary dependency.
\item Discard the remaining signed parities, pad each neighborhood
to $k$ distinct inputs, and compute the proper Fourier layers.
\item Construct the degree-zero and degree-one certificates directly.
For even higher degrees, repeatedly remove a dense window after a
failed averaged trace test. For odd degrees, construct equalized
pair matrices and repeatedly remove the dense forests supplied by
failed window LPs. Retain every resulting local and remainder
certificate.
\item Enumerate a single sufficiently small-biased space on the
original output signs. Return a vector passing the summed
certificate test in Eqn.~\eqref{eq:joint-layer-budget}, extended
arbitrarily on discarded coordinates.
\end{enumerate}
Every branch uses exact rational tests. The algorithm does not need
to certify the local capacity condition on all small windows at a
stopping point: the averaged trace test itself is the remainder
certificate used in the proof.

% End avoidance


% Begin hierarchy
\section{The Hierarchy Tradeoff}\label{sec:hierarchy}

We now allow the matrix level to grow with $n$ and prove
\Cref{thm:hierarchy-avoid}. For $d\geq2$ and $1\leq\ell\leq n$, define
\begin{equation}\label{hierarchy:scale}
 \mu_d(n,\ell)=n\left(\frac n\ell\right)^{d/2-1}
              =\ell\left(\frac n\ell\right)^{d/2}.
\end{equation}
The role of the level is to increase the average matrix degree produced
by each original term. The local capacity and cancellation parameters
remain constants depending only on $d$ and the desired margin.

\begin{proposition}\label{hierarchy:signing}
Fix $d\geq2$ and rational $0<\epsilon<1$. The certificate conclusion
of \Cref{thm:signing} holds under the density hypothesis
$M\geq C_{d,\epsilon}\mu_d(n,\ell)$, with construction, evaluation,
and inverse bias bounded by
\begin{equation}\label{hierarchy:cost}
 \poly(\text{input})
 \left(\frac{\mathrm e n}{\ell}\right)^{O_{d,\epsilon}(\ell)}.
\end{equation}
The expectation guarantee continues to hold for every final distribution
with the specified full bias.
\end{proposition}

\subsection{The central range}

Assume first that $2d\leq\ell$ and $n\geq100d\ell$, and write
$v=n/\ell$ and $N=\binom n\ell$. For $1\leq q\leq d$, the exact
slice ratio is
\begin{equation}\label{hierarchy:slice-ratio}
 \frac{\binom{n-2q}{\ell-q}}{\binom n\ell}
 =\frac{(\ell)_q(n-\ell)_q}{(n)_{2q}}
 =\Theta_d\bigl((\ell/n)^q\bigr),
\end{equation}
where $(a)_q=a(a-1)\cdots(a-q+1)$.
The constants are independent of $\ell$: the numerator factors are
bounded below by $(\ell/2)^q(n/2)^q$ and above by $\ell^q n^q$,
and $(n/2)^{2q}\leq(n)_{2q}\leq n^{2q}$.

\paragraph{Even degree.}
Suppose $d=2r$. Use the padded level-$\ell$ Kikuchi graph: a clause
$F_i$ joins $S$ and $S\triangle F_i$ whenever $\abs{S\cap F_i}=r$.
The number of ordered instances per original clause is now
\[
 c_{d,\ell}=\binom dr\binom{n-d}{\ell-r}.
\]
Thus $\overline d=c_{d,\ell}W_R/N$, and the comparison identities in
Eqn.~\eqref{even:comparison} hold with $c_{d,\ell}$ in place of $c$.
Each label still determines at most one neighbor of a given row, so
the weighted Hall allocation and orientation proof is unchanged.
Its incoming capacity is $(D+1)\ell$.

By Eqn.~\eqref{hierarchy:slice-ratio},
$N/c_{d,\ell}=\Theta_d(v^{d/2})$. The sufficient weighted density
condition from Eqn.~\eqref{even:density} therefore has the form
\[
 W\geq C'_{d,\epsilon}\ell\frac{N}{c_{d,\ell}}.
\]
It is implied by $W\geq C_{d,\epsilon}\mu_d(n,\ell)$ for a sufficiently
large constant. With this replacement,
the deterministic decomposition, dense-block moments, and spectral
certificate from \Cref{sec:even} apply verbatim.
To convert weighted density into label-count density, return the
constant certificate $W$ if $W\leq\epsilon M/2$.
Otherwise $W>\epsilon M/2$, and a sufficiently large constant in
$M\geq C_{d,\epsilon}\mu_d(n,\ell)$ ensures the weighted hypothesis.
The weighted certificate then has expectation at most
$(41/128)\epsilon W\leq\epsilon M/2$.

\paragraph{Odd degree.}
Suppose $d\geq3$ is odd. Retain the constants $B,R$ and the
cancellation budgets of \Cref{sec:odd}, and replace the bundle
thresholds by
\begin{equation}\label{hierarchy:bundles}
 h_q=\left\lceil BR^q
          \max\{1,v^{q-d/2}\}\right\rceil,
 \qquad 0\leq q\leq d-1.
\end{equation}
The packing proof uses only the exhaustion of earlier stages and
$h_{q-1}/h_q\leq2/R$, so its spread and multiplicity conclusions
continue to hold. The number of unbundled labels is at most
\begin{equation}\label{hierarchy:unbundled}
 \frac{nh_{d-1}}d
 =O_{d,\epsilon}\bigl(nv^{d/2-1}\bigr)
 =O_{d,\epsilon}\bigl(\mu_d(n,\ell)\bigr).
\end{equation}

Fix a positive residual size $q$ and put
$a=\lfloor q/2\rfloor$, $b=\lceil q/2\rceil$.
Each disjoint original pair has
\[
 T_q=\begin{cases}
 \frac12\binom qa^2\binom{n-2q}{\ell-q},&q\text{ even},\\
 \binom qa^2\binom{n-2q}{\ell-q},&q\text{ odd}
 \end{cases}
\]
balanced undirected instances. Hence $T_q/N=\Theta_d(v^{-q})$.
We check that the contention constant remains independent of $\ell$.
For a competing residual whose intersection with the fixed partner has
size $t<a$, the eligibility probability is $O_d(v^{-(a-t)})$.
For each fixed intersection, spread bounds the number of competitors by
$h_{q-t}$, with $h_q$ used at $t=0$. Their contribution is at most
\[
 h_{q-t}O_d(v^{-(a-t)})
 \leq O_d(BR^q)
       \bigl(v^{-(a-t)}+v^{b-d/2}\bigr)
 \leq O_d(BR^q).
\]
The last inequality uses $b\leq(d-1)/2<d/2$.
If $t\geq a$, then $q-t\leq b<d/2$, so $h_{q-t}$ itself is a
constant depending only on $d,B,R$. This includes $t=q$ and therefore
also covers repeated residuals. Summing over the constantly many
intersections and four ports gives a uniform contention bound.
Exact pair equalization consequently supplies a common pair weight
$\beta$ with
\begin{equation}\label{hierarchy:pair-weight}
 \frac\beta N\geq c_{d,\epsilon}v^{-q}.
\end{equation}

The good-window orientation still has incoming capacity $(D+1)\ell$.
Choose the degree stopping threshold $d_*$ to be the same sufficiently
large constant multiple of $\ell$ used in the trace argument.
At a level containing $M_q$ original labels, the direct low-degree
remainder is bounded by
\begin{align}
 \frac{Nd_*}{\beta M_qh_q}
 &\leq\frac{C_{d,\epsilon}\ell v^q}
                {M_q\max\{1,v^{q-d/2}\}}\notag\\
 &\leq C_{d,\epsilon}\frac{\ell v^{d/2}}{M_q}
  =C_{d,\epsilon}\frac{\mu_d(n,\ell)}{M_q}.
 \label{hierarchy:remainder}
\end{align}
Thus a sufficiently large threshold
$M_q\geq A_{d,\epsilon}\mu_d(n,\ell)$ pays the same fixed squared
error budget as before. Discarding smaller positive levels costs only
$O_{d,\epsilon}(\mu_d(n,\ell))$ labels because there are at most
$d-1$ such levels. Diagonal, overlap, extracted-forest, and scalar-block
budgets are unchanged. The common-sign certificate assembly of
\Cref{sec:odd} now gives the asserted expectation after increasing
$C_{d,\epsilon}$.

\paragraph{Cost of the central constructions.}
For both parities, choose $s=O_{d,\epsilon}(\log N)$ with the same
trace thresholds as in the fixed-level proof. Since
\[
 \ell\log(n/\ell)\leq\log N
 \leq\ell\log(\mathrm e n/\ell),
\]
and $n/\ell\geq100d$, the memory windows have
$L=O_d(\ell+ds)=O_{d,\epsilon}(\log N)$ vertices.
Their assignment enumerations cost $N^{O_{d,\epsilon}(1)}$.
The trace threshold has inverse polynomial in $N$, as do the
heavy-memory extraction threshold and the required full bias:
the normalized traces have Fourier mass at most $N$, and the local
moment potentials have mass $2^{O_{d,\epsilon}(L)}$.
All graph dimensions, flow or linear programs, matrix powers, and
local certificates therefore cost polynomial time in the explicit
input and $N$. Frozen pair weights and rational arithmetic have
polynomial bit length in these same parameters. This proves
Eqn.~\eqref{hierarchy:cost} in the central range.

\subsection{Boundary levels and scalar derandomization}

If $\ell<2d$ and $n\geq200d^2$, run the central construction at
$\ell'=2d$. Since $d/2-1\geq0$,
$\mu_d(n,\ell')\leq\mu_d(n,\ell)$, so its density requirement is
met. Its polynomial running time fits Eqn.~\eqref{hierarchy:cost}.

If $n<100d\ell$, use the whole-input scalar certificate of
\Cref{lem:scalar}, with $u=n$ and $e=\epsilon/2$.
The required $M\geq16e^{-2}(n+1)$ follows from
$M\geq C_{d,\epsilon}\mu_d(n,\ell)$ because $\mu_d(n,\ell)\geq n$.
Its expectation is at most $\epsilon M/2$ under the specified bias,
whose inverse is $2^{O_{d,\epsilon}(n)}$.
Construction and evaluation have the same exponential dependence on
$n$, and
\[
 2^{O_{d,\epsilon}(n)}
 \leq(\mathrm e n/\ell)^{O_{d,\epsilon}(\ell)}
\]
because $n<100d\ell$ and $\log(\mathrm e n/\ell)\geq1$.
Any case left by these two rules has $n<200d^2$ and is handled by
the same scalar construction with constant overhead depending on
$d,\epsilon$. This proves \Cref{hierarchy:signing}.

For completeness, a scalar signing can also be found directly by
conditional expectation of its even-moment potential. This procedure
enumerates input assignments, and never enumerates $2^M$ sign vectors.
Fix an input $x$, put $c_i=a_i\chi_{F_i}(x)$ and $p=2(n+1)$, and
precompute the suffix moments
\[
 H_j(r)=\E\left(\sum_{i=j}^M c_i\xi_i\right)^r,
 \qquad 0\leq r\leq p,
\]
where the $\xi_i$ are independent uniform signs. Initialize
$H_{M+1}(0)=1$ and $H_{M+1}(r)=0$ for $r>0$, and use
\begin{equation}\label{hierarchy:moment-recurrence}
 H_j(r)=\sum_{\substack{0\leq h\leq r\\h\ \mathrm{even}}}
         \binom rh c_j^h H_{j+1}(r-h).
\end{equation}
If the chosen prefix contributes $z$ at $x$, fixing the next sign to
$b_j=\pm1$ gives conditional $p$th moment
\[
 \sum_{r=0}^p\binom pr(z\pm c_j)^{p-r}H_{j+1}(r).
\]
Sum these expressions over all inputs with the normalization in
\Cref{lem:scalar}, and choose a sign whose conditional expectation
does not exceed the current one. The uniform starting expectation
supplies the desired certificate bound at termination.
The recurrence and all conditional comparisons take
$\poly(\text{input})2^{O(n)}$ bit operations. The same computation
minimizes a sum of finitely many scalar certificates on common signs,
since conditional expectations add.

\subsection{The avoidance tradeoff}

As in \Cref{sec:avoidance}, $n<k$ is handled by evaluating at most
$2^k$ inputs, so assume $n\geq k$ and apply that section's Fourier
reduction. If more than $n$
outputs are signed parities, its linear-algebra branch finishes the
problem. Otherwise, remove those outputs and retain $M\geq m-n$
nonparity outputs. Their proper Fourier parts have degree at most
$k-1$, with at most $J=2^k-1$ layers and Fourier gap
$\delta=2^{1-k}$.

For every layer of degree $2\leq d\leq k-1$, use
\Cref{hierarchy:signing} with $\epsilon=\delta/(4J)$.
Since $n/\ell\geq1$,
\[
 \mu_d(n,\ell)\leq\mu_{k-1}(n,\ell)
 =n(n/\ell)^{(k-3)/2}.
\]
Degree-zero and degree-one layers use the exact certificates from
\Cref{sec:avoidance}. Their expectations are bounded by
$\sqrt M+\sqrt\eta M$ and $\sqrt{nM}+\sqrt\eta M$, respectively,
so a sufficiently large constant multiple of $n$ and sufficiently
small fixed bias give the same margin. Since
$\mu_{k-1}(n,\ell)\geq n$, one constant $C_k$ absorbs these
requirements and the discarded parity outputs.

Choose a common full bias no larger than any layer requires and
enumerate one explicit biased space. The sum of the certificates
has expectation at most $\delta M/8$, so some enumerated vector
has sum strictly below $\delta M$. The contradiction in
\Cref{sec:avoidance} then proves that this vector, with arbitrary
values in the discarded positions, is outside the range.
There are only $O_k(1)$ layers; taking the smallest of their bias
parameters and adding their construction and evaluation costs
preserves Eqn.~\eqref{hierarchy:cost}. This is precisely
Eqn.~\eqref{eq:hierarchy-time} under
Eqn.~\eqref{eq:hierarchy-stretch}, proving
\Cref{thm:hierarchy-avoid}.

% End hierarchy


% Begin applications
\section{One-Sided Distinguishers and Data Structure Lower Bounds}
\label{sec:applications}

Finding one missing string need not distinguish a generator from uniform:
the test that accepts just that string has acceptance probability $2^{-m}$.
The expectation guarantee in \Cref{thm:signing} supplies the stronger
conclusion needed here. We first retain that guarantee through the
Fourier reduction and then apply it to cryptography and data structures.
These applications follow the connections developed by Korten, Pitassi,
and Impagliazzo~\cite{KPI25} and Guruswami, Lyu, and Yuan~\cite{GLY26}.
All locality and accuracy parameters in this section are fixed.

\subsection{A Dense, Efficiently Recognizable Subset of the Complement}

\begin{theorem}[One-sided range test]\label{thm:onesided}
For each fixed $k\ge3$, there are constants $C_k,B_k>0$ and a
deterministic polynomial-time algorithm that, given an explicit
$k$-local map $f:\{0,1\}^n\to\{0,1\}^m$ with
$m\ge C_kn^{(k-1)/2}$, constructs a polynomial-size Boolean circuit
$D_f:\{0,1\}^m\to\{0,1\}$ satisfying
\begin{align}
 D_f(f(x))&=0 &&\text{for every }x,\label{app:sound}\\
 \Pr_{y\sim\U_m}[D_f(y)=1]&\ge\tfrac12.\label{app:uniform}
\end{align}
Moreover, for every fully $\eta$-biased distribution $\mathcal Q$ with
$\eta\le(n+m+2)^{-B_k}$,
\begin{equation}\label{app:biased}
 \Pr_{y\sim\mathcal Q}[D_f(y)=1]\ge\tfrac13.
\end{equation}
The circuit $D_f$ is constructed without access to $\mathcal Q$.
\end{theorem}

\begin{proof}
Assume first $n\ge k$. If there are more than $n$ signed-parity
coordinates, find a nonempty binary dependency among them as in
\Cref{sec:avoidance}. Let $D_f$ accept precisely when the corresponding
output parity violates the identity obeyed by $f$. This proves
Eqn.~\eqref{app:sound}. Its uniform acceptance probability is $1/2$;
under $\mathcal Q$ it is at least $(1-\eta)/2$.

Otherwise retain the $M\ge m-n$ nonparity coordinates. Use exactly the
certificates constructed in \Cref{sec:avoidance}, and put
\[
 T_f(y)=\sum_{A\subsetneq[k]}\mathcal C_A(y),
 \qquad \delta=2^{1-k}.
\]
Here and below each certificate reads only the retained coordinates
of $y$. Every image point satisfies $T_f(f(x))\ge\delta M$, while
Eqn.~\eqref{eq:joint-layer-budget} gives
\[
 \E_{\mathcal Q}T_f\le\delta M/8
\]
whenever the bias is small enough for every layer. Define
$D_f(y)=\mathbf1\{T_f(y)<\delta M\}$. Soundness follows from the
first inequality, and Markov's inequality gives acceptance probability
at least $7/8$ under both uniform and sufficiently small-biased inputs.
Construction, exact rational comparison, and Boolean circuit simulation
of the evaluation algorithm have polynomial bit complexity.

There is a single exponent $B_k$ that suffices for all instances.
Indeed, their truth-table Fourier coefficients have $O_k(1)$ bit length,
and the inverse biases in the fixed-parameter constructions are bounded
by a polynomial in the explicit representation length
$O_k(m\log(n+2)+\log(n+2))$. Increasing $B_k$ absorbs constants and
ensures also $\eta\le1/3$.

If $n<k$, enumerate the image and let $D_f$ recognize its complement.
Uniform acceptance is at least $1-2^{n-m}\ge1/2$. For a fully biased
distribution, the character expansion of a singleton gives
$\mathcal Q(y)\le2^{-m}+\eta$. Thus
$\mathcal Q(\Range(f))\le1/2+2^k\eta\le2/3$ after a further constant
increase of $B_k$. This computation has polynomial cost at fixed $k$.
\end{proof}

In the nonparity branch the test can certify a constant Hamming gap.
If $T_f(y)\le\delta M/2$, then
\[
 \sum_{i=1}^M y_i f_i(x)
 \le(1-\delta)M+T_f(y)\le(1-\delta/2)M.
\]
Consequently every image differs from $y$ in at least $\delta M/4$
retained positions. This event has probability at least $3/4$.
There is no corresponding constant-distance assertion for the
single-parity test in the other branch.

\begin{corollary}[Any fixed coefficient]\label{cor:prefactor}
Fix $k\ge4$ and rational $\rho>0$. The conclusions of
\Cref{thm:onesided}, with constants depending on $k,\rho$, hold whenever
$m>n$ and $m\ge\rho n^{(k-1)/2}$. Construction and evaluation take time
\[
 (n+m+2)^{O_k(1+\rho^{-2/(k-3)})}.
\]
\end{corollary}
\begin{proof}
Choose the constant level
$\ell_0=\max\{1,\lceil(C_k/\rho)^{2/(k-3)}\rceil\}$ in
\Cref{thm:hierarchy-avoid}. For $n\ge\ell_0$, its density requirement
is at most $\rho n^{(k-1)/2}$. The hierarchy proof preserves the
common-bias expectation guarantee before choosing a signing, so the
proof of \Cref{thm:onesided} applies. For $n<\ell_0$, use exhaustive
image evaluation and the same singleton bound, increasing the
bias exponent by a constant depending on $k,\rho$.
The time bound follows from $\ell_0=O_k(1+\rho^{-2/(k-3)})$.
\end{proof}

\subsection{Consequences for Local Demi-Bits and Goldreich Generators}

We use the convention of Ren, Wang, and Zhong~\cite[Definition~2.1]{RWZ26}:
an $(s,\varepsilon)$-secure demi-bits generator has no nondeterministic
circuit of size at most $s$ that rejects its entire image and accepts
at least an $\varepsilon$ fraction of uniform strings.

\begin{corollary}[Breaking local demi-bits]\label{cor:demibits}
For every fixed $k\ge3$, every polynomial-output-length family of
$k$-local maps satisfying $m(n)\ge C_kn^{(k-1)/2}$ admits
polynomial-size deterministic distinguishers that reject the entire
image and accept at least half of all uniform strings. In particular,
such a family is not secure as a demi-bits generator against
$\NP/\poly$. For $k\ge4$, the same conclusion holds with any fixed
positive coefficient in place of $C_k$, provided $m>n$.
\end{corollary}
\begin{proof}
Apply \Cref{thm:onesided,cor:prefactor}. A deterministic circuit is
a special case of a nondeterministic circuit. For a nonuniform
generator family, hardwire its polynomial-length local description
and the constructed test. If the local descriptions are uniformly
computable, the adversary is uniform as well.
\end{proof}

The attack is one-sided and therefore also distinguishes the output
distribution of an ordinary local pseudorandom generator from uniform.
If $m\ge rC_kn^{(k-1)/2}$, apply the test to $r$ disjoint output
blocks and accept when any test accepts. Uniform blocks are independent,
so this raises the acceptance probability to at least $1-2^{-r}$;
every image point is still rejected. No independence among generator
outputs is used.

For Goldreich's generator, the predicate and the ordered support
hypergraph may be arbitrary. For fixed locality $k\ge4$, a family
with output length $m=n^c$ is therefore excluded as a secure
demi-bits generator whenever $c\ge(k-1)/2$. In particular, for
the five-variable predicate
\[
 P(x)=x_1+x_2+x_3+x_4x_5\quad\text{over }\F_2,
\]
the conclusion applies at $m=\rho n^2$ for every fixed $\rho>0$.
This predicate is discussed as a candidate in~\cite[Appendix~B.3]{RWZ26}.
The result does not exclude its instantiation at $m=n^{1+\varepsilon}$
for $0<\varepsilon<1$. At locality three, the conclusion requires a
sufficiently large linear-stretch constant, not every $m>n$ or every
specified constant multiple of $n$.

Constant degree over $\F_2$ is different from constant locality.
For example, a quadratic output may involve all input variables.
Accordingly, \Cref{cor:demibits} does not refute the general
constant-degree demi-bits assumption in~\cite[Assumption~2.3]{RWZ26}.
It restricts its local candidate instantiations in the parameter
range just stated. Li, Ren, and Zhong~\cite[Corollary~4.6]{LRZ26}
already deduce a nondeterministic demi-bits attack at the GLY
threshold. Our endpoint removes its logarithmic loss, and
\Cref{thm:onesided} supplies a deterministic one-sided test directly.

\paragraph{The projection barrier to further improvements.}
Li, Ren, and Zhong~\cite[Corollary~4.5]{LRZ26} show that, for a
universal constant $c$, a secure $k$-local demi-bits generator of
output length $c s(n)$ rules out even an errorless nondeterministic
polynomial-time solver for $k$-local AVOID at output length $s(n)>n$.
Thus the existence of secure local demi-bits at a smaller stretch
would obstruct a corresponding improvement of our algorithm.
This is a conditional barrier, not an unconditional optimality
claim for the exponent $(k-1)/2$.

\paragraph{Proof complexity.}
There is also a direct consequence at the level of general
propositional proof systems. For each fixed $k$, extend a standard
complete proof system by a rule that accepts a range-exclusion
tautology for $(f,y)$ when the promise of \Cref{thm:onesided} holds
and $D_f(y)=1$. The rule is sound and polynomial-time checkable.
It gives polynomial-size proofs for at least half of the strings
$y$ at the stated stretch. This defines a Cook--Reckhow proof system;
it is not a claim of short proofs in Resolution, Polynomial Calculus,
or a fixed-degree sum-of-squares system. Establishing any such claim
requires a separate simulation of the certificate construction.

The projection argument of Li, Ren, and
Zhong~\cite[Theorem~3.2 and Section~4]{LRZ26} strengthens the
uniform-output guarantee. We record its quantitative consequence;
the amplification mechanism is theirs.

\begin{corollary}[Short certificates for almost all candidate outputs]
\label{cor:projection-proofs}
Fix $k\ge3$, and let $s,m$ be integers with
\[
 n<s\le m,\qquad s\ge C_kn^{(k-1)/2}.
\]
There is a polynomial-time verifier $V_k(f,y,I)$, uniform in the
local description of $f:\{0,1\}^n\to\{0,1\}^m$, whose witness
$I$ is an $s$-element subset of $[m]$, such that
\begin{enumerate}
 \item $V_k(f,y,I)=1$ implies $y\notin\operatorname{Range}(f)$;
 \item the number of strings having no accepting witness is at most
 \begin{equation}\label{app:projection-exceptions}
  \sum_{j=0}^{s-1}\binom{m}{j}.
 \end{equation}
\end{enumerate}
In particular, if $m\ge4s$, the verifier accepts, for some witness,
at least a $1-2^{-\gamma m}$ fraction of uniform strings, where
$\gamma=1-h_2(1/4)>0$ and
$h_2(u)=-u\log_2u-(1-u)\log_2(1-u)$.
For each fixed $k$, one Cook--Reckhow proof system gives
polynomial-size proofs of the corresponding range-exclusion
tautologies for all these accepted strings.
\end{corollary}
\begin{proof}
Let $A_k$ be the deterministic algorithm of \Cref{thm:avoid}.
The verifier first rejects malformed local descriptions and witnesses
that violate the stated length and density conditions.
Given $I$, form the projected map $f|_I$ in increasing coordinate
order, compute $z_I=A_k(f|_I)$, and accept precisely when
$y|_I=z_I$. Projection preserves locality and the density promise
holds at output length $s$. Thus $z_I$ is outside the range of
$f|_I$, proving soundness. An $m$-bit incidence vector encodes $I$,
and verification takes polynomial time for fixed $k$.

Let $U$ be the set of strings with no accepting witness. For every
$s$-element set $I$, the projection $U|_I$ omits $z_I$.
Consequently $U$ shatters no set of $s$ coordinates. The
Sauer--Shelah bound, recalled in~\cite[Lemma~2.10]{LRZ26}, gives
Eqn.~\eqref{app:projection-exceptions}. If $s\le m/4$, the
binomial entropy bound gives
\[
 2^{-m}|U|
 \le 2^{-m}\sum_{j=0}^{s-1}\binom{m}{j}
 \le 2^{-m(1-h_2(1/4))}.
\]
Finally, extend a fixed complete Cook--Reckhow system by a proof
format containing $(f,y,I)$, checked by $V_k$, and producing the
canonical range-exclusion tautology for $(f,y)$. The verifier is
sound and polynomial-time, and the base system preserves
completeness. Its proof lengths are polynomial in the local
description and output length.
\end{proof}

More generally, $s\le\beta m$ for fixed $0<\beta<1/2$ gives
exceptional fraction at most $2^{-(1-h_2(\beta))m}$.
At $s=\lceil C_kn^{(k-1)/2}\rceil$, increasing $C_k$ if needed
to ensure $s>n$, the condition $m\ge4s$ changes only the constant
in the main stretch bound. This conclusion uses the AVOID
algorithm alone. Finding a witness $I$ for a given $y$ is not
shown to be deterministic polynomial time, and the exponentially
small exceptional fraction is asserted for uniform $y$, not for
every small-biased distribution. The deterministic test and
small-bias guarantee of \Cref{thm:onesided} remain separate inputs
to the data structure applications. Likewise, the proof system
constructed here does not imply short proofs in an independently
specified weak system. Transferring the result to such a system
requires short proofs for the projected instances and a suitable
closure under adding axioms; see~\cite[Section~4.3]{LRZ26}.

These are average-case guarantees over candidate outputs of each fixed
map. They do not constitute a worst-case-to-average-case reduction
over a fixed distribution of circuit instances. The conditional
average-case hardness result in~\cite[Theorem~4.4]{LRZ26} instead
starts from a demi-bits generator and samples its output projections.

\subsection{Universal Families and Nonadaptive Bit Probes}

\begin{corollary}[An oblivious family of candidate nonoutputs]
\label{cor:universal}
For fixed $k\ge3$ and lengths $n,m$ satisfying the density promise,
there is an explicitly enumerable multiset $\mathcal H_{n,m}$ of
$(n+m)^{O_k(1)}$ strings such that, for every $k$-local map
$f:\{0,1\}^n\to\{0,1\}^m$, at least one third of the elements of
$\mathcal H_{n,m}$, counted with multiplicity, are accepted by $D_f$.
The multiset is chosen before $f$ is specified.
\end{corollary}
\begin{proof}
Use \Cref{lem:small-bias} with $h=m$ and
$\eta=(n+m+2)^{-B_k}$, and apply Eqn.~\eqref{app:biased}.
\end{proof}

The quantifiers are
$\exists\mathcal H\ \forall f\ \exists y\in\mathcal H$,
not $\exists y\ \forall f$. They are exactly what the data structure
application needs. A fixed query algorithm must represent every row
of the query problem, although each row has its own memory encoding.

A static query problem is a Boolean matrix $F:[R]\times[m]\to\{0,1\}$.
An encoding assigns each row a memory string of $S$ bits. A nonadaptive
$t$-probe query algorithm reads at most $t$ memory locations, determined
by the query index, and returns the requested entry. No computational
bound is placed on encoding or decoding.

\begin{lemma}[Strongly explicit rows]\label{lem:explicit-rows}
For every fixed integer $B\ge1$ and $m\ge1$, a fully
$(3m+3)^{-B}$-biased distribution on $m$ bits can be indexed by
$O_B(\log(m+2))$ uniform seed bits. Any specified output coordinate
is computable deterministically in time $\poly_B(\log(m+2))$ from
the seed, the coordinate index, and $m$.
\end{lemma}
\begin{proof}
Put $\eta=(3m+3)^{-B}$ and choose $r$ to be the least power of two
at least $\lceil\log_2\max\{2,m/\eta\}\rceil$. Use the construction
in \Cref{lem:small-bias} over $\F_{2^r}$, indexed by
$(a,z)\in\F_{2^r}^2$. Its bias is at most $(m-1)/2^r\le\eta$,
and the seed length is $2r=O_B(\log(m+2))$.

For completeness, this field has a deterministic quadratic-tower
construction of cost polynomial in $r$. Begin with $K=\F_2$ and
$c=1$, whose absolute trace is one. Replace $K$ by
$K[u]/(u^2+u+c)$. The polynomial is irreducible because
$\operatorname{Tr}_{K/\F_2}(c)=1$: the image of $v\mapsto v^2+v$
is exactly the trace-zero subspace. In the new field, $cu$ has
absolute trace one, since its relative trace to $K$ is $c$.
Use $cu$ in the next extension. Repeating until dimension $r$
gives arithmetic in a specified binary basis in time polynomial
in $r$. Computing $z^{j-1}$ and its trace then has cost polynomial
in $r+\log m$. This avoids any need for an exponential search for
an irreducible polynomial when evaluating a query.
\end{proof}

\begin{theorem}[Logarithm-free nonadaptive space bound]
\label{thm:bitprobe}
For every fixed $t\ge3$, there is $B_t>0$ such that the following holds.
Let the rows of $F:[R]\times[m]\to\{0,1\}$ support a fully
$(3m+3)^{-B_t}$-biased distribution. Every exact nonadaptive
$t$-bit-probe representation of $F$ requires
\begin{equation}\label{app:bitprobe}
 S\ge c_t m^{2/(t-1)}
\end{equation}
for a constant $c_t>0$. There are explicit such matrices with
$R=m^{O_t(1)}$ and entries computable in time $\poly_t(\log(m+2))$.
\end{theorem}
\begin{proof}
If $S\ge m$, Eqn.~\eqref{app:bitprobe} is immediate after taking
$c_t\le1$. Otherwise a fixed decoder induces the $t$-local map
\[
 C:\{0,1\}^S\longrightarrow\{0,1\}^m,
 \qquad C(z)_j=\text{the answer to query }j\text{ on memory }z.
\]
Correctness for every row implies
$\supp(\mathcal Q)\subseteq\Range(C)$. If
$m\ge C_tS^{(t-1)/2}$, \Cref{thm:onesided} contradicts this
containment: the chosen bias is at most $(S+m+2)^{-B_t}$ and
the test accepts a positive fraction of $\mathcal Q$ outside the range.
Hence $m<C_tS^{(t-1)/2}$, which gives the bound. Zero-bit memories
are impossible for a distribution with bias less than one.

For explicitness, use \Cref{lem:explicit-rows} on $m$ bits at
the stated bias. Its sample points are the rows, with multiplicity.
The seed length is $O_t(\log(m+2))$ and its coordinate evaluation
has the asserted running time.
\end{proof}

In particular, $t=3$ gives $S=\Omega(m)$, while $t=4$ gives
$S=\Omega(m^{2/3})$ and $t=5$ gives $S=\Omega(m^{1/2})$.
The row index has only $O_t(\log m)$ bits. The three-probe statement
is therefore an explicit linear-space lower bound for a problem
whose data can be specified with logarithmically many bits.
The earlier reduction in~\cite[Theorem~10]{GLY26} gives the same
power with a logarithmic loss. Our statement removes that loss for
a sufficiently small polynomial bias; we do not preserve their
specific bias exponent.

\subsection{Adaptive Queries and Constant-Size Words}

We give a separate consequence for adaptive cell probes. It uses the
full degree-$t$ expansion, rather than treating a depth-$t$ decision
tree as a junta on $2^t-1$ bits. This distinction avoids an exponential
loss in the exponent of the storage bound.

\begin{theorem}[Remote points for bounded-depth word queries]
\label{thm:wordremote}
Fix integers $t\ge2,w\ge1$ and rational $0<\varepsilon<1/2$.
There are constants $C,B>0$ depending on these parameters such that,
for integers $S,m\ge1$, if each coordinate of
$f:(\{0,1\}^w)^S\to\{\pm1\}^m$ is computed by an adaptive
decision tree making at most $t$ word queries and
$m\ge CS^{t/2}$, a deterministic polynomial-time test certifies
\[
 \min_z\Delta(f(z),y)\ge\tfrac12-\varepsilon
\]
with probability at least $3/4$ under every fully
$(S+m+2)^{-B}$-biased distribution on $y$.
The test takes the explicit query trees as input and never certifies
a false assertion.
\end{theorem}
\begin{proof}
For every cell $v$ and nonzero $a\in\F_2^w$, introduce the formal
sign variable $X_{v,a}=(-1)^{\langle a,z_v\rangle}$.
There are $N=(2^w-1)S$ such features. Delete repeated queries to an
already known cell along any path of a decision tree. The indicator
of a leaf is a product of at most $t$ word-value indicators, and
\[
 \mathbf1\{z_v=c\}=2^{-w}
       \sum_{a\in\F_2^w}(-1)^{\langle a,c\rangle}X_{v,a},
 \qquad X_{v,0}=1.
\]
The resulting Fourier expansion has degree at most $t$ in the
formal features and at most $J=2^{2wt}$ nonzero monomials per output.
After collecting terms, every coefficient has magnitude at most one:
these are Fourier coefficients of a Boolean function on the original
independent uniform cells. Each monomial uses at most one feature
from any cell.

Order the monomials separately by degree in each output, pad lists
with zero coefficients, and form at most $(t+1)J$ systems, with one
original label per output in every system. For degrees at least two,
apply \Cref{thm:signing} on the $N$ formal variables. Treating the
features as independent enlarges the domain and hence preserves
the upper-bound direction. Degrees zero and one use
\Cref{lem:elementary-degrees}. Pad the formal variable set to size
at least $t$ if necessary; this changes only fixed constants.

Choose each accuracy parameter as
$\varepsilon/((t+1)J)$. The density hypothesis, with sufficiently
large $C$, satisfies every layer's requirement. Their sum $T(y)$
is nonnegative, efficiently evaluable, and obeys
\[
 \sup_z\left|\sum_i y_if_i(z)\right|\le T(y),
 \qquad \E T\le\varepsilon m/2.
\]
All inverse biases are uniformly polynomial in $S+m$ for fixed
$t,w,\varepsilon$. Markov's inequality gives
$T(y)\le2\varepsilon m$ with probability at least $3/4$.
The identity
$\Delta(f(z),y)=(1-m^{-1}\sum_i y_if_i(z))/2$
completes the proof.
\end{proof}

\begin{corollary}[Adaptive constant-word space bound]
\label{cor:wordprobe}
For fixed $t\ge2,w\ge1$, there is an explicit query problem with
$m^{O_{t,w}(1)}$ rows such that every exact adaptive $t$-probe
representation using $w$-bit cells requires
\[
 S=\Omega_{t,w}(m^{2/t}).
\]
The same bound holds if each row need only be represented within
relative Hamming error $1/4$.
\end{corollary}
\begin{proof}
Take $\varepsilon=1/8$ in \Cref{thm:wordremote} and choose a common
bias $(3m+3)^{-B}$ for all $S<m$. Use \Cref{lem:explicit-rows}.
Its explicit small-biased rows
cannot all lie within distance $1/4$ of the decoder's range,
because a positive fraction have distance at least $3/8$.
The case $S\ge m$ is immediate. A zero-cell decoder is constant;
its average distance from a small-biased row is at least
$(1-\eta)/2>1/4$, so it cannot provide the stated approximation.
\end{proof}

The word size is fixed: this corollary makes no assertion of uniform
dependence on a growing $w=\Theta(\log m)$. Also, when $w=1$, the
specialized adaptive bit-probe analysis of~\cite[Section~5.2]{GLY26}
has stronger powers of $m$ than $2/t$. Thus \Cref{cor:wordprobe}
is a logarithm-free baseline and a robust approximation statement,
not a replacement for those stronger adaptive bit-probe results.
All data structure assertions above concern exact deterministic
decoding, or the explicitly stated per-row Hamming approximation;
bounded-error randomized decoding requires a separate argument.

\section{What the Sharp Moore Bound Supplies Directly}
\label{sec:moore-applications}

The sharp density scale in our construction comes from the recent
proofs of the hypergraph Moore bound~\cite{BKNPW26,SH26Moore}.
Bandeira et al.~\cite[Theorem~1.2, version~2]{BKNPW26} prove that,
for every fixed uniformity $d\ge3$, density
$M\ge c_d n(n/\ell)^{d/2-1}$ forces an even cover of size
at most $C_d\ell\log n$. Their second version
covers odd as well as even $d$.
Schmidhuber and Hastings additionally state the
refined cover-length bound $O(\ell\log(en/\ell))$ with absolute
constants. Their random-instance companion~\cite{SH26Kikuchi}
does not by itself supply our arbitrary-support, common-sign
certificate theorem. This section records a direct consequence of
the extremal result and separates it from the constructive work.

\begin{proposition}[Many disjoint short dependencies]
\label{prop:moore-packing}
Fix $d\ge3$ and $1\le\ell\le n$. There are constants $A,C>0$
such that every labeled $d$-uniform hypergraph with
\[
 M\ge2R,\qquad
 R=\left\lceil Cn(n/\ell)^{d/2-1}\right\rceil
\]
has a collection of pairwise label-disjoint nonempty even covers,
each of size at most
$L=\max\{2,\lceil A\ell\log(en/\ell)\rceil\}$,
whose union contains at least $M/2$ labels.
\end{proposition}
\begin{proof}
While at least $R$ labels remain, find an even cover and delete it.
Two repeated supports immediately give a cover of size two;
otherwise apply the sharp Moore bound to the simple residual
hypergraph. Choose $C$ large enough for its strict density
inequality. Fewer than $R$ labels remain at termination, so at
least $M/2$ have been removed.
\end{proof}

\begin{corollary}[A weak refutation margin from short covers]
\label{cor:moore-weak}
Under the hypotheses of \Cref{prop:moore-packing}, consider the
signed equations $\chi_{F_i}(x)=b_i$. Under every fully
$\eta$-biased distribution on the $M$ labels, with probability
at least $1-8L/M-4\eta$, every assignment violates at least
$M/(8L)$ equations. Equivalently, the maximum normalized
correlation is at most $1-1/(4L)$.
\end{corollary}
\begin{proof}
Fix the packing before drawing the signs, and let it have $q$ covers
$E_1,\ldots,E_q$. Since their union has at least $M/2$ labels,
$q\ge M/(2L)$. Put $Z_j=\prod_{i\in E_j}b_i$. Every $Z_j=-1$
certifies an inconsistent subsystem: multiplying its equations
has left-hand side one. The covers are label-disjoint, so every
$Z_jZ_h$ with $j\ne h$ is a nonempty parity. Therefore
\[
 \E\left(\sum_{j=1}^q Z_j\right)^2
 \le q+\eta q(q-1).
\]
Markov's inequality bounds the probability that
$\sum_jZ_j>q/2$ by $4/q+4\eta\le8L/M+4\eta$.
On the complementary event, at least $q/4$ covers are inconsistent.
Each forces a violation on a different label, proving the claim.
\end{proof}

The covers provide polynomial-size verifiable witnesses even when
this existence proof does not find them in polynomial time at a
growing level. The weak margin is $\Theta(1/L)$, not a fixed
constant. Passing from these witnesses to efficiently constructed
strong cancellation for arbitrary supports is precisely why the
memory extraction, pair equalization, and forest decomposition in
the preceding sections are needed. Neither the short-cover theorem
alone nor the present fixed-accuracy analysis proves a smaller
polynomial-time stretch exponent or automatically removes every
logarithmic factor from the specialized adaptive-query bounds.

% End applications


\section*{Acknowledgments}
This draft was prepared during the author's visit to the Simons Institute
for the Theory of Computing for the program Pseudorandomness and
High-Dimensional Expansion. The author thanks the institute for its
hospitality.

\paragraph{AI Disclosure.}
ChatGPT/Codex was used extensively in developing the arguments, checking
proof details, reviewing the literature, and preparing this draft.
Computational checks of the forest and memory-extraction components
supplement the mathematical arguments. These checks do not implement
the full asymptotic algorithm, and no formal verification is claimed.

\newcommand{\etalchar}[1]{$^{#1}$}
\begin{thebibliography}{BKNN{\etalchar{+}}26}

\bibitem[AGHP92]{AGHP92}
Noga Alon, Oded Goldreich, Johan H{\aa}stad, and Ren{\'e} Peralta.
\newblock Simple construction of almost {$k$}-wise independent random
  variables.
\newblock {\em Random Structures \& Algorithms}, 3(3):289--304, 1992.

\bibitem[BKNN{\etalchar{+}}26]{BKNPW26}
Afonso~S. Bandeira, Dmitriy Kunisky, Petar Nizi{\'c}-Nikolac, Lucas Pesenti,
  and Robert Wang.
\newblock The hypergraph {Moore} bound.
\newblock {\em CoRR}, abs/2607.14068, 2026.
\newblock arXiv version 2.

\bibitem[GGNS23]{GGNS23}
Karthik Gajulapalli, Alexander Golovnev, Satyajeet Nagargoje, and Sidhant
  Saraogi.
\newblock Range avoidance for constant depth circuits: Hardness and algorithms.
\newblock In {\em APPROX/RANDOM 2023}, volume 275 of {\em Leibniz International
  Proceedings in Informatics (LIPIcs)}, pages 65:1--65:18. Schloss
  Dagstuhl--Leibniz-Zentrum f\"ur Informatik, 2023.

\bibitem[GKM22]{GKM22}
Venkatesan Guruswami, Pravesh~K. Kothari, and Peter Manohar.
\newblock Algorithms and certificates for {Boolean CSP} refutation: Smoothed is
  no harder than random.
\newblock In {\em 54th Annual ACM SIGACT Symposium on Theory of Computing (STOC
  2022)}, pages 678--689. Association for Computing Machinery, 2022.

\bibitem[GLW22]{GLW22}
Venkatesan Guruswami, Xin Lyu, and Xiuhan Wang.
\newblock Range avoidance for low-depth circuits and connections to
  pseudorandomness.
\newblock In {\em APPROX/RANDOM 2022}, volume 245 of {\em Leibniz International
  Proceedings in Informatics (LIPIcs)}, pages 20:1--20:21. Schloss
  Dagstuhl--Leibniz-Zentrum f\"ur Informatik, 2022.

\bibitem[GLY25]{GLY26}
Venkatesan Guruswami, Xin Lyu, and Weiqiang Yuan.
\newblock Cell-probe lower bounds via semi-random {CSP} refutation: Simplified
  and the odd-locality case.
\newblock {\em CoRR}, abs/2507.22265, 2025.
\newblock arXiv version 1.

\bibitem[HKM23]{HKM23}
Jun-Ting Hsieh, Pravesh~K. Kothari, and Sidhanth Mohanty.
\newblock A simple and sharper proof of the hypergraph {Moore} bound.
\newblock In {\em 2023 Annual ACM-SIAM Symposium on Discrete Algorithms
  (SODA)}, pages 2324--2344. Society for Industrial and Applied Mathematics,
  2023.

\bibitem[KPI25]{KPI25}
Oliver Korten, Toniann Pitassi, and Russell Impagliazzo.
\newblock Stronger cell probe lower bounds via local {PRGs}.
\newblock {\em Electronic Colloquium on Computational Complexity}, TR25-030,
  2025.

\bibitem[LRZ26]{LRZ26}
Xin Li, Hanlin Ren, and Yan Zhong.
\newblock Many proof complexity generators inside one demi-bits generator.
\newblock {\em CoRR}, abs/2609.23228, 2026.
\newblock arXiv version 1; also ECCC TR26-196.

\bibitem[NN93]{NN93}
Joseph Naor and Moni Naor.
\newblock Small-bias probability spaces: Efficient constructions and
  applications.
\newblock {\em SIAM Journal on Computing}, 22(4):838--856, 1993.

\bibitem[RWZ26]{RWZ26}
Hanlin Ren, Yichuan Wang, and Yan Zhong.
\newblock Hardness of range avoidance and proof complexity generators from
  demi-bits.
\newblock {\em CoRR}, abs/2511.14061, 2026.
\newblock arXiv version 2.

\bibitem[SH26a]{SH26Kikuchi}
Alexander Schmidhuber and Matthew~B. Hastings.
\newblock The {Kikuchi} hierarchy is sharp for {$k$XOR}.
\newblock {\em CoRR}, abs/2607.29672, 2026.
\newblock arXiv version 1.

\bibitem[SH26b]{SH26Moore}
Alexander Schmidhuber and Matthew~B. Hastings.
\newblock A spectral proof of the hypergraph {Moore} bound.
\newblock {\em CoRR}, abs/2607.26028, 2026.
\newblock arXiv version 1.

\end{thebibliography}

\end{document}
